\chapter{非完整系统的运动规划与控制：从 Brockett 障碍到链式型与平坦化}
\label{ch:nonholonomic}

% ════════════════════════════════════════════════════════════════
%  控制部 P4。三书有机集大成（非翻译、重述）：
%   - MLS（Murray-Li-Sastry）Ch7/Ch8：几何控制理论深度（李括号、可控性、正弦转向、链式型）。
%   - Siciliano Ch11：轮式建模与工程控制（跟踪误差、近似线性化、非线性 Lyapunov、极坐标镇定）。
%   - Spong Ch14：Brockett 障碍的精确陈述、动态反馈线性化、开/闭环对照。
%  分工纪律：底盘 Pfaffian 矩阵的逐一推导、Frobenius/Chow/LARC 的判据细节，已由轮式运动学
%   章 ch:mm-wheeled 给足，本章只做最小复述并 \cref 过去；本章的真增量是「规划与控制」半边：
%   Brockett 障碍、链式型、微分平坦、正弦/多项式转向、轨迹跟踪 vs 定点镇定。
%  记号归一：q=(x,y,θ) in SE(2)，控制向量场 g_i，仿射系统 q̇=Σ u_i g_i = G(q)u，与 ch:mm-wheeled 一致。
% ════════════════════════════════════════════════════════════════

\begin{practice}[前置自测]
开始之前先自测；卡住之处正是本章要打通的。
\begin{enumerate}
  \item 差速车“瞬时不能横移”，为何仍能平行泊车挪进侧方车位？哪个数学对象精确刻画了这“凭空多出的侧向运动”（\cref{ch:mm-wheeled}）？
  \item 线性系统 $\dot{\mathbf{x}}=\mathbf{A}\mathbf{x}+\mathbf{B}\mathbf{u}$ 若可控，能否用光滑状态反馈 $\mathbf{u}=-\mathbf{K}\mathbf{x}$ 把它指数镇定到原点（\cref{sec:ctrl-controllability}）？换成独轮车，同样的“可控就能光滑镇定”还成立吗？
  \item 欠驱动系统（\cref{ch:underactuated}）受限于“执行器少于自由度”；非完整系统受限于什么？二者是同一回事吗？
  \item 什么是微分平坦？若一个系统平坦，规划一条可行轨迹为何会变得几乎“免费”？
\end{enumerate}
\end{practice}

\section*{本章目标}
\addcontentsline{toc}{section}{本章目标}
读完本章，你应当能够：
\begin{enumerate}
  \item \textbf{定位}非完整约束在“受限系统”谱系中的位置——它与欠驱动（执行器不足）、与完整约束（位形被困在子流形）的本质区别，是\emph{速度约束不可积}；
  \item \textbf{讲透} Brockett 必要条件：为什么\emph{光滑时不变状态反馈无法把非完整系统渐近镇定到一点}，以及这一“硬障碍”如何把全部镇定方案逼向\emph{时变 / 不连续 / 切换}三条路；
  \item \textbf{掌握}链式型（chained form）与微分平坦（differential flatness）：把独轮车、自行车的运动学化成规范形 $\dot z_1=v_1,\dot z_2=v_2,\dot z_3=z_2 v_1,\dots$，并由平坦输出代数地回算全状态与全输入；
  \item \textbf{会用}转向（steering）：正弦输入（Murray--Sastry）、多项式/平坦轨迹两类开环转向，把系统从 $\mathbf{q}_i$ 驱到 $\mathbf{q}_f$；
  \item \textbf{区分并实现}两类闭环任务：\emph{轨迹跟踪}（近似线性化、非线性 Lyapunov、I/O 线性化、动态反馈线性化——皆可镇定）与\emph{定点姿态镇定}（时变 / 不连续 / $\sigma$ 过程——结构性地更难），并理解“跟踪可、镇定难”的根由正是 Brockett 障碍。
\end{enumerate}

\subsection*{本章知识导航}
\begin{center}
\begin{tikzpicture}[
  node distance=4.5mm and 7mm, >=Stealth,
  blk/.style={draw, rounded corners, align=center, font=\small, inner sep=3pt, fill=main!6, draw=main!60!black},
  grp/.style={draw, rounded corners, align=center, font=\small, inner sep=3pt, fill=second!6, draw=second!60!black},
  hot/.style={draw, rounded corners, align=center, font=\small, inner sep=3pt, fill=third!10, draw=third!65!black},
  ln/.style={->, thick, gray!60!black}]
\node[blk] (con) {约束与可控性\\（复述 \cref{ch:mm-wheeled}）};
\node[hot, right=of con] (bro) {\textbf{Brockett 障碍}\\光滑镇定不可能};
\node[grp, right=of bro] (cf) {链式型 + 平坦\\规范形 / 平坦输出};
\node[grp, below=8mm of con] (steer) {开环转向\\正弦 / 多项式};
\node[blk, below=8mm of bro] (trk) {轨迹跟踪\\近似/非线性/动态 FBL};
\node[hot, below=8mm of cf] (reg) {定点镇定\\时变 / 不连续 / $\sigma$};
\draw[ln] (con)--(bro); \draw[ln] (bro)--(cf);
\draw[ln] (cf.south)-- (steer.north -| cf.south) ;
\draw[ln] (cf)--(trk); \draw[ln] (cf)--(reg);
\draw[ln] (bro.south) -- (trk.north);
\draw[ln] (bro.south) -- (reg.north);
\end{tikzpicture}
\end{center}

\noindent\textbf{推荐路径}：\cref{sec:nh-place,sec:nh-recap} 把前置（受限系统谱系、约束/可控性）一句话收齐，已读过 \cref{ch:mm-wheeled} 者可速览；\cref{sec:nh-brockett} 是\emph{全章枢纽}，把“为何非完整系统难镇定”讲到根上；\cref{sec:nh-chained,sec:nh-flat} 是规划的代数引擎（链式型 + 平坦）；\cref{sec:nh-steer} 给开环转向；\cref{sec:nh-track,sec:nh-reg} 给两类闭环任务，是工程落点。

\subsection*{前置知识桥接}
\begin{itemize}
  \item \textbf{受限系统谱系}：欠驱动（\cref{ch:underactuated}）与本章并列——前者“输入维数 $<$ 自由度”，后者“速度约束不可积”；二者都使系统\emph{不可线性化为全可控}，但机理不同（\cref{sec:nh-place}）。
  \item \textbf{约束分类与可控性}：完整/非完整的形式定义（\cref{def:gm-constraint}）、Pfaffian 形式 $\mathbf{A}(\mathbf{q})\dot{\mathbf{q}}=\mathbf{0}$（\cref{eq:cd-pfaffian}）、Frobenius 可积判据（\cref{thm:cd-frobenius}）、四种底盘的 Pfaffian 矩阵与 Chow--Rashevsky/LARC（\cref{ch:mm-wheeled}）——本章直接调用，不重证。
  \item \textbf{线性可控性}：\cref{sec:ctrl-controllability} 的可控性是本章\emph{非线性几何可控性}的特例与对照；本章反复用“线性化是否可控”作为分水岭。
  \item \textbf{李括号与几何力学}：向量场对易子 $[\mathbf{g}_1,\mathbf{g}_2]$ 的定义（\cref{ch:lie}）与漂移自由仿射系统（\cref{ch:geometric-mechanics}）。
  \item \textbf{Lyapunov 与 Barbalat}：定点镇定与跟踪的收敛性证明全靠 \cref{sec:ctrl-lyapunov} 的 Lyapunov 方法，外加时变系统专用的 Barbalat 引理。
\end{itemize}

\subsection*{如果跳过本章会怎样}
\begin{itemize}
  \item 你会以为“差速车可控（\cref{ch:mm-wheeled} 已证），那写个 $\mathbf{u}=-\mathbf{K}\mathbf{q}$ 反馈就能把它停到指定位姿”——一跑就发现机器人在目标点附近\emph{打转不收敛}；不懂 Brockett 障碍，就找不到病根（\cref{sec:nh-brockett}）。
  \item 你会把轨迹跟踪和定点泊车当成同一个问题用同一个控制器解——结果跟踪控制器（要求参考轨迹“持续激励”）一旦目标静止就失效（\cref{sec:nh-reg}）。
  \item 规划时你会去硬解一个带非完整微分约束的边值问题，费力且易失败；而\emph{平坦}本可让你在平坦输出空间里随便插值、再代数回算（\cref{sec:nh-flat,sec:nh-steer}），轻松保证可行。
\end{itemize}

% ════════════════════════════════════════════════════════════════
\section{定位：非完整 vs 欠驱动 vs 完整}
\label{sec:nh-place}

“受限”在机器人里有三种相貌不同的来源，初学者极易混为一谈。把它们一次摆清，本章的目标就立住了。

\begin{definition}[三类“受限”的辨析]\label{def:nh-taxonomy}
设系统位形 $\mathbf{q}\in\mathcal{Q}$（$\dim\mathcal{Q}=n$），输入 $\mathbf{u}\in\mathbb{R}^m$。
\begin{itemize}
  \item \textbf{完整约束（holonomic）}：可写成位形的代数等式 $\mathbf{h}(\mathbf{q})=\mathbf{0}$。它\emph{削减位形空间维数}——把 $\mathbf{q}$ 困在 $\mathcal{Q}$ 的子流形上（\cref{def:gm-constraint}）。例：刚性连杆 $\lVert\mathbf{p}_1-\mathbf{p}_2\rVert=\ell$。
  \item \textbf{非完整约束（nonholonomic）}：对广义速度的线性约束 $\mathbf{A}(\mathbf{q})\dot{\mathbf{q}}=\mathbf{0}$（Pfaffian），且\emph{不可积}为任何 $\mathbf{h}(\mathbf{q})=\mathbf{0}$。它\emph{削减瞬时速度自由度、却不削减可达位形}——门后的房间一间不少，只是每步运动方向被上了锁。例：车轮纯滚动不横滑。
  \item \textbf{欠驱动（underactuated）}：$m<n$，独立执行器少于自由度（\cref{ch:underactuated}）。它\emph{削减瞬时可施控方向}，约束来自\emph{执行器}而非\emph{运动学}。例：Acrobot（被动第一关节）、平面 PVTOL。
\end{itemize}
\end{definition}

三者的共性是都让系统\emph{不能在每一时刻往任意方向走}，因而其线性化往往\emph{不完全可控}，标准线性反馈设计失灵。但机理与后果各异，下表把它们摆在一起。

\begin{table}[H]\centering\small
\caption{完整 / 非完整 / 欠驱动：削减的对象与典型后果}
\label{tab:nh-taxonomy}
\begin{tabular}{llll}
\toprule
类型 & 约束形式 & 削减的是 & 可达位形受限？ \\
\midrule
完整 & $\mathbf{h}(\mathbf{q})=\mathbf{0}$ & 位形空间维数 & 是（困在子流形） \\
非完整 & $\mathbf{A}(\mathbf{q})\dot{\mathbf{q}}=\mathbf{0}$ 不可积 & 瞬时速度自由度 & \textbf{否}（仍可达全空间，但要绕路） \\
欠驱动 & $m<n$ 执行器 & 瞬时可施控方向 & 视具体系统（可能含漂移项，更复杂） \\
\bottomrule
\end{tabular}
\end{table}

\begin{insight}[非完整：唯一“瞬时受限、整体自由”者]\label{ins:nh-essence}
非完整最反直觉之处，恰在于\cref{tab:nh-taxonomy} 中间一行那个加粗的“否”。完整约束把可达集压成低维曲面；非完整约束\emph{一寸位形都没削掉}，它只是规定了到达任一位形\emph{必须沿哪些方向、按什么顺序}走。差速车能到平面上任意 $(x,y,\theta)$——靠“前进—转弯—前进”的机动序列，其可达性由 Chow 定理担保（\cref{ch:mm-wheeled} 已证 LARC 成立）。把“瞬时不能横移”误读成“到不了侧方”，是本主题的头号误区。
\end{insight}

\begin{pitfall}[误区：非完整 $=$ 欠驱动]
独轮车 $\mathbf{q}=(x,y,\theta)$、输入 $\mathbf{u}=(v,\omega)$，恰好 $m=2<n=3$，看着像欠驱动。但它“受限”的根源是\emph{车轮不横滑}这条速度约束，不是“少装了一个电机”。区别的判据：给独轮车再加一个能直接产生横向速度的执行器（变成全向底盘），约束消失、系统变完整可控；而 Acrobot 哪怕你换更强的电机，被动关节的\emph{欠驱动}也还在。本章只处理\emph{无漂移}（driftless）的纯运动学非完整系统 $\dot{\mathbf{q}}=\sum_i u_i\mathbf{g}_i(\mathbf{q})$；含惯性漂移的二阶非完整与欠驱动力学，属 \cref{ch:underactuated} 与本部动力学诸章。
\end{pitfall}

% ════════════════════════════════════════════════════════════════
\section{约束与可控性：最小必要复述}
\label{sec:nh-recap}

轮式底盘的 Pfaffian 矩阵推导、Frobenius 与 Chow 判据，\cref{ch:mm-wheeled} 已逐一给足。本节只把后文要反复用的\emph{结论}收成一页，不重证；不熟者请回看该章。

\paragraph{从约束到无漂移控制系统。}
非完整约束统一写成 Pfaffian 齐次形式 $\mathbf{A}(\mathbf{q})\dot{\mathbf{q}}=\mathbf{0}$。其零空间 $\mathcal{D}(\mathbf{q})=\ker\mathbf{A}(\mathbf{q})$ 给出一组控制向量场 $\{\mathbf{g}_1,\dots,\mathbf{g}_m\}$（$m=n-\operatorname{rank}\mathbf{A}$），系统随之写成\emph{漂移自由仿射系统}的标准形
\begin{equation}
  \dot{\mathbf{q}}=\sum_{i=1}^{m}u_i\,\mathbf{g}_i(\mathbf{q})=\mathbf{G}(\mathbf{q})\,\mathbf{u},
  \qquad \mathbf{G}(\mathbf{q})=[\mathbf{g}_1(\mathbf{q})\ \cdots\ \mathbf{g}_m(\mathbf{q})].
  \label{eq:nh-affine}
\end{equation}
“漂移自由”意指无与输入无关的 $\mathbf{f}_0(\mathbf{q})$ 项：输入归零则机器人静止。本章一切结论都挂靠在 \cref{eq:nh-affine}。

\begin{example}[独轮车 / 差速 / 自行车的标准形（取自 \cref{ch:mm-wheeled}）]\label{ex:nh-models}
取 $\mathbf{q}=(x,y,\theta)^\top\in\SE(2)$，输入为前向速度 $v$、转向角速度 $\omega$。
\begin{itemize}
  \item \textbf{独轮车 / 差速}（差速底盘经轮速 $\to(v,\omega)$ 映射后与独轮车同形）：唯一约束“横向不滑” $\dot x\sin\theta-\dot y\cos\theta=0$，标准形
  \begin{equation}
    \begin{bmatrix}\dot x\\\dot y\\\dot\theta\end{bmatrix}
    =\underbrace{\begin{bmatrix}\cos\theta\\\sin\theta\\0\end{bmatrix}}_{\mathbf{g}_1}v
    +\underbrace{\begin{bmatrix}0\\0\\1\end{bmatrix}}_{\mathbf{g}_2}\omega .
    \label{eq:nh-unicycle}
  \end{equation}
  \item \textbf{自行车 / Ackermann}（后轮驱动，加转向角 $\phi$，$\mathbf{q}=(x,y,\theta,\phi)^\top$，轴距 $\ell$）：
  \begin{equation}
    \begin{bmatrix}\dot x\\\dot y\\\dot\theta\\\dot\phi\end{bmatrix}
    =\begin{bmatrix}\cos\theta\\\sin\theta\\\tan\phi/\ell\\0\end{bmatrix}v
    +\begin{bmatrix}0\\0\\0\\1\end{bmatrix}\omega ,
    \label{eq:nh-bicycle}
  \end{equation}
  此处 $v$ 为后轮线速度、$\omega$ 为转向角速度 $\dot\phi$。当 $\phi\to\pm\pi/2$ 时 $\tan\phi$ 发散，对应转向打死、模型失效（最小转弯半径约束）。
\end{itemize}
\end{example}

\paragraph{可控性：李括号生成“侧向”。}
对 \cref{eq:nh-affine}，李括号 $[\mathbf{g}_i,\mathbf{g}_j]=\frac{\partial\mathbf{g}_j}{\partial\mathbf{q}}\mathbf{g}_i-\frac{\partial\mathbf{g}_i}{\partial\mathbf{q}}\mathbf{g}_j$ 给出“原向量场所不能直达、却可由机动序列凑出”的新运动方向。\textbf{Chow--Rashevsky 定理}（李代数秩条件 LARC）：若 $\{\mathbf{g}_i\}$ 及其各阶李括号张成的李代数在每点维数等于 $n$，则系统\emph{小时间局部可控}（STLC）。\cref{ch:mm-wheeled} 已算出独轮车的 $[\mathbf{g}_1,\mathbf{g}_2]=(\sin\theta,-\cos\theta,0)^\top$ 恰是\emph{横向}方向，与 $\mathbf{g}_1,\mathbf{g}_2$ 并列张满 $\SE(2)$——这正是“泊车为何可行”的数学解释。

\begin{insight}[同一个李括号，可积性与可控性的两面（呼应 \cref{ch:mm-wheeled}）]\label{ins:nh-twoface}
Frobenius 与 Chow 用的是\emph{同一个}李括号 $[\mathbf{g}_1,\mathbf{g}_2]$：Frobenius 问它\emph{在不在}速度分布 $\mathcal{D}$ 内（在 $\Rightarrow$ 约束可积、运动被困）；Chow 问它\emph{补没补齐}全空间（补齐 $\Rightarrow$ 系统可控）。对独轮车，$[\mathbf{g}_1,\mathbf{g}_2]$ 逸出 $2$ 维分布（故非完整）、又恰补成 $3$ 维（故可控）。本章下一步要追问的是：可控\emph{未必}意味着“能被光滑反馈镇定到一点”——这正是 Brockett 障碍登场处。
\end{insight}

% ════════════════════════════════════════════════════════════════
\section{Brockett 障碍：光滑镇定的不可能性}
\label{sec:nh-brockett}

这是全章枢纽。对\emph{线性可控}系统，可控即可用光滑反馈 $\mathbf{u}=-\mathbf{K}\mathbf{x}$ 指数镇定到原点（\cref{sec:ctrl-controllability}）。许多人想当然地把这条结论搬到非完整系统：“Chow 定理证了差速车可控，那它也该能被某个 $\mathbf{u}(\mathbf{q})$ 光滑镇定到目标位姿。”\textbf{这是错的}，而且错得很深——它是非完整控制最根本的障碍。

\begin{theorem}[Brockett 必要条件\cite{brockett1983asymptotic}]\label{thm:nh-brockett}
设非线性系统 $\dot{\mathbf{x}}=\mathbf{f}(\mathbf{x},\mathbf{u})$，$\mathbf{f}$ 连续可微、$\mathbf{f}(\mathbf{0},\mathbf{0})=\mathbf{0}$。若存在连续可微（光滑）的时不变状态反馈 $\mathbf{u}=\mathbf{u}(\mathbf{x})$ 使原点 $\mathbf{x}=\mathbf{0}$ 渐近稳定，则映射 $\mathbf{f}$ 必须\emph{局部到上}（局部开映射）：$\mathbf{f}$ 的像 $\{\mathbf{f}(\mathbf{x},\mathbf{u}):(\mathbf{x},\mathbf{u})\in\mathcal{N}\}$ 包含原点的一个邻域。
\end{theorem}

\begin{remark}[“像含原点邻域”是什么意思]\label{rem:nh-brockett-meaning}
直观地：要光滑镇定，系统必须有能力在原点附近\emph{往每个方向产生净位移}——给定任意微小目标方向 $\boldsymbol{\epsilon}$，都得能找到 $(\mathbf{x},\mathbf{u})$ 使 $\mathbf{f}(\mathbf{x},\mathbf{u})=\boldsymbol{\epsilon}$。这是个\emph{瞬时}的、代数的覆盖条件，与 Chow 的\emph{积分}可达性截然不同：Chow 只要“多步机动后能到”，Brockett 要“一步就能往该方向使劲”。非完整系统恰恰瞬时缺了某些方向（横向），于是过不了 Brockett 这关。
\end{remark}

\begin{example}[Brockett 系统：链式型的最小反例]\label{ex:nh-brockett-fail}
取 $(2,3)$ 链式系统（独轮车经坐标变换后的规范形，见 \cref{sec:nh-chained}）
\begin{equation}
  \dot x_1=u_1,\qquad \dot x_2=u_2,\qquad \dot x_3=u_1 x_2 .
  \label{eq:nh-brockett-sys}
\end{equation}
检验 Brockett 条件：能否对任意小 $a\neq0$ 解出 $(\mathbf{x},\mathbf{u})$ 使 $\mathbf{f}(\mathbf{x},\mathbf{u})=(0,0,a)^\top$？需 $u_1=0$（第一式）、$u_2=0$（第二式），但那样第三式给 $\dot x_3=u_1 x_2=0\neq a$。\emph{矛盾}。故 $\mathbf{f}$ 的像在 $x_3$ 轴方向上漏掉了原点邻域——\cref{thm:nh-brockett} 的必要条件不满足，\textbf{不存在}光滑时不变反馈把 \cref{eq:nh-brockett-sys} 渐近镇定到原点。独轮车、差速、自行车、拖挂链\emph{全部}属此类。
\end{example}

\begin{insight}[障碍的几何根：李括号方向“瞬时使不上劲”]\label{ins:nh-brockett-geo}
\cref{ex:nh-brockett-fail} 里漏掉的 $x_3$ 方向，正是 \cref{ins:nh-twoface} 里那个\emph{靠李括号 $[\mathbf{g}_1,\mathbf{g}_2]$ 凑出来}的方向。系统能到达 $x_3$ 目标（Chow 担保），但只能靠“先动 $u_1$ 再动 $u_2$ 再回来”的\emph{二阶}效应净挪一段——任何\emph{瞬时}的输入都在该方向产生零速度。光滑反馈要求镇定速度随状态连续地、单调地指向目标；而非完整系统在李括号方向上\emph{瞬时速度恒为零}，光滑反馈便无从在该方向“连续地往回拉”，只能靠不光滑的来回抖动或随时间编排的机动。一句话：\emph{可控（能绕路到）$\neq$ 光滑可镇定（能被连续反馈直接拉到）}。
\end{insight}

\begin{theorem}[非完整镇定三出路（Brockett 障碍的后果）]\label{thm:nh-three-ways}
由于绝大多数非完整系统不满足 \cref{thm:nh-brockett}，要把它们镇定到一点，必须放弃“光滑 + 时不变 + 纯状态反馈”三者中的至少一个。可行的三类方案是：
\begin{enumerate}
  \item \textbf{时变反馈} $\mathbf{u}=\mathbf{u}(\mathbf{q},t)$：引入显含时间（常取周期）的项，绕过时不变性限制；
  \item \textbf{不连续 / 切换反馈} $\mathbf{u}=\mathbf{u}(\mathbf{q})$（分段、含开关面）：放弃光滑性，如滑模、$\sigma$ 过程、极坐标律（在原点不连续）；
  \item \textbf{开环（前馈）}：不闭环到目标，按规划好的输入序列（正弦 / 多项式）把系统驱过去（\cref{sec:nh-steer}），对扰动不鲁棒但简单。
\end{enumerate}
\end{theorem}

\begin{pitfall}[误区：Brockett 也禁止跟踪]
Brockett 障碍\emph{只}针对“镇定到一个\emph{固定点}”。若目标是跟踪一条\emph{持续运动}的参考轨迹（参考速度不趋零、“持续激励”），障碍\emph{不}适用——此时沿轨迹线性化得到的时变系统可控，光滑（时变）反馈即可指数跟踪（\cref{sec:nh-track}）。正因如此，本章把跟踪与镇定分成两节：跟踪是“好做”的（\cref{sec:nh-track}），定点镇定才撞上 Brockett 障碍、必须动用 \cref{thm:nh-three-ways} 的特设手段（\cref{sec:nh-reg}）。这条“跟踪可、定点镇定难”的分界，是非完整控制的工程主旋律。
\end{pitfall}

% ════════════════════════════════════════════════════════════════
\section{链式型：非完整系统的规范形}
\label{sec:nh-chained}

要系统化地规划与控制，先把五花八门的底盘运动学化成\emph{统一的规范形}。对两输入系统（独轮车、自行车、拖挂链），这个规范形就是\textbf{链式型}（chained form），由 Murray 与 Sastry 引入\cite{murray1993sinusoids}。它的好处是：可控性一目了然、可闭式积分、平坦输出现成。

\begin{definition}[$(2,n)$ 链式型]\label{def:nh-chained}
$(2,n)$ 链式型是带两个输入的无漂移系统
\begin{equation}
  \dot z_1=v_1,\quad \dot z_2=v_2,\quad \dot z_3=z_2 v_1,\quad \dot z_4=z_3 v_1,\ \dots,\ \dot z_n=z_{n-1}v_1 .
  \label{eq:nh-chained}
\end{equation}
其结构是“一条链”：$v_1$ 驱动 $z_1$，并经 $z_2,z_3,\dots$ 逐级\emph{积分}向下传递；$v_2$ 只直接驱动 $z_2$。
\end{definition}

\begin{proposition}[链式型必可控，且非完整度 $\kappa=n-1$]\label{prop:nh-chained-ctrl}
记 $\boldsymbol{\gamma}_1=(1,0,z_2,z_3,\dots,z_{n-1})^\top$、$\boldsymbol{\gamma}_2=(0,1,0,\dots,0)^\top$ 为 \cref{eq:nh-chained} 的两控制向量场。用重复李括号记号 $\operatorname{ad}_{\boldsymbol{\gamma}_1}^k\boldsymbol{\gamma}_2=[\boldsymbol{\gamma}_1,\operatorname{ad}_{\boldsymbol{\gamma}_1}^{k-1}\boldsymbol{\gamma}_2]$，可逐级算得 $\operatorname{ad}_{\boldsymbol{\gamma}_1}^k\boldsymbol{\gamma}_2$ 的第 $(k+2)$ 个分量为 $(-1)^k$、其余为零。于是可达分布
\begin{equation}
  \Delta_{\mathcal{A}}=\operatorname{span}\{\boldsymbol{\gamma}_1,\boldsymbol{\gamma}_2,\operatorname{ad}_{\boldsymbol{\gamma}_1}\boldsymbol{\gamma}_2,\dots,\operatorname{ad}_{\boldsymbol{\gamma}_1}^{n-2}\boldsymbol{\gamma}_2\}
  \label{eq:nh-chained-acc}
\end{equation}
处处满秩 $n$，LARC 成立、系统可控（\cref{ins:nh-twoface}）；且要张满需 $n-2$ 阶李括号，故\emph{非完整度} $\kappa=n-1$。链长越长（拖挂越多），非完整度越高，规划越难（Ball--Box 代价 $O(\varepsilon^{1/\kappa})$，\cref{ch:mm-wheeled}）。
\end{proposition}

\paragraph{独轮车 $\to(2,3)$ 链式型。}
对独轮车 \cref{eq:nh-unicycle}，用坐标与输入变换
\begin{equation}
  z_1=\theta,\quad z_2=x\cos\theta+y\sin\theta,\quad z_3=x\sin\theta-y\cos\theta;
  \qquad v=v_2+z_3 v_1,\quad \omega=v_1,
  \label{eq:nh-uni2chain}
\end{equation}
即得 $(2,3)$ 链式型 $\dot z_1=v_1,\ \dot z_2=v_2,\ \dot z_3=z_2 v_1$。其中 $z_1$ 就是朝向 $\theta$，而 $(z_2,z_3)$ 是机器人位置在\emph{随车转动}的坐标系下的表示（$z_2$ 轴沿车体纵轴）。

\begin{proof}[变换的核验思路]
对 \cref{eq:nh-uni2chain} 第二式求时间导数：$\dot z_2=\dot x\cos\theta+\dot y\sin\theta+(-x\sin\theta+y\cos\theta)\dot\theta$。代入 \cref{eq:nh-unicycle} 的 $\dot x=v\cos\theta,\dot y=v\sin\theta,\dot\theta=\omega$，前两项合为 $v$、括号项为 $-z_3\omega$，故 $\dot z_2=v-z_3\omega$。再代输入变换 $v=v_2+z_3v_1,\omega=v_1$ 得 $\dot z_2=v_2$。同法验 $\dot z_1=\dot\theta=\omega=v_1$、$\dot z_3=z_2 v_1$。逐式吻合 $\square$
\end{proof}

\paragraph{自行车 $\to(2,4)$ 链式型。}
后轮驱动自行车 \cref{eq:nh-bicycle} 经坐标 $z_1=x,\ z_2=\frac{1}{\ell}\sec^3\theta\tan\phi,\ z_3=\tan\theta,\ z_4=y$ 与相应输入变换，化为 $(2,4)$ 链式型 $\dot z_1=v_1,\dot z_2=v_2,\dot z_3=z_2 v_1,\dot z_4=z_3 v_1$。变换在 $\cos\theta=0$（$\theta=\pm k\pi/2$）处奇异，故等价性须避开车头垂直于 $x$ 轴的位形。

\begin{note}[何时一定能化成链式型]\label{note:nh-chainable}
存在把一般两输入无漂移系统 \cref{eq:nh-affine} 化成链式型的\emph{充分必要}条件（涉及一族分布 $\Delta_0,\Delta_1,\Delta_2$ 的对合性与一个满足 $\mathrm{d}h_1\cdot\Delta_1=0,\ \mathrm{d}h_1\cdot\mathbf{g}_1=1$ 的标量函数 $h_1$ 的存在性，构造性地给出变换\cite{murray1993sinusoids}）。一个干净的充分结论：维数 $n\le4$ 的两输入无漂移系统\emph{总能}化成链式型——独轮车（$n=3$）、自行车（$n=4$）都在内。更长的拖挂链需逐情形验证。
\end{note}

% ════════════════════════════════════════════════════════════════
\section{微分平坦：可行轨迹几乎免费}
\label{sec:nh-flat}

链式型的真正威力，要等到看清它\emph{平坦}才显出来。微分平坦把“规划一条满足非完整约束的可行轨迹”这个看似要解微分约束边值问题的难题，降格成“在低维输出空间里随便插一条光滑曲线”。

\begin{definition}[微分平坦与平坦输出]\label{def:nh-flat}
系统 $\dot{\mathbf{x}}=\mathbf{f}(\mathbf{x})+\mathbf{G}(\mathbf{x})\mathbf{u}$（$\mathbf{x}\in\mathbb{R}^n,\mathbf{u}\in\mathbb{R}^m$）称为\emph{微分平坦}的，若存在一组与输入同维的输出 $\mathbf{y}\in\mathbb{R}^m$（\emph{平坦输出}），使状态与输入都能写成 $\mathbf{y}$ 及其有限阶时间导数的\emph{代数}函数：
\begin{equation}
  \mathbf{x}=\boldsymbol{\Phi}_x\!\left(\mathbf{y},\dot{\mathbf{y}},\dots,\mathbf{y}^{(r)}\right),
  \qquad
  \mathbf{u}=\boldsymbol{\Phi}_u\!\left(\mathbf{y},\dot{\mathbf{y}},\dots,\mathbf{y}^{(r+1)}\right).
  \label{eq:nh-flat}
\end{equation}
\end{definition}

\begin{theorem}[无漂移系统：平坦 $\Leftrightarrow$ 可化链式型]\label{thm:nh-flat-chained}
对无漂移系统（如轮式底盘的运动学模型），微分平坦与可化为链式型\emph{等价}。特别地，$(2,n)$ 链式型 \cref{eq:nh-chained} 的平坦输出是 $z_1$ 与 $z_n$（链头与链尾）：由 $z_1(t),z_n(t)$ 及其导数可代数地回算全部中间状态与两输入。
\end{theorem}

\begin{example}[独轮车的笛卡尔坐标即平坦输出]\label{ex:nh-flat-unicycle}
取平坦输出 $\mathbf{y}=(x,y)$（机器人位置）。给定任意光滑笛卡尔路径 $(x(t),y(t))$（$\dot x,\dot y$ 不同时为零），由 \cref{eq:nh-unicycle} 代数地回算其余量：
\begin{align}
  \theta &= \operatorname{Atan2}(\dot y,\dot x),
  & v &= \pm\sqrt{\dot x^2+\dot y^2},
  & \omega &= \dot\theta=\frac{\ddot y\,\dot x-\ddot x\,\dot y}{\dot x^2+\dot y^2}.
  \label{eq:nh-flat-recon}
\end{align}
朝向是速度矢量的方位角（$v$ 取正/负对应前进/倒车），转向角速度是路径曲率乘以速率。\emph{结论}：随手画一条 $C^2$ 曲线 $(x(t),y(t))$，\cref{eq:nh-flat-recon} 当场吐出一条\emph{自动满足非完整约束}的完整状态—输入轨迹。规划的难度被“代数化”掉了。
\end{example}

\begin{insight}[平坦为何让规划“免费”]\label{ins:nh-flat-why}
不平坦时，规划可行轨迹要在 $n$ 维状态空间里解一个受非完整微分约束的两点边值问题——又难又易失败。平坦把这件事翻译到 $m$ 维平坦输出空间（独轮车 $m=2$）：那里\emph{没有微分约束}，任何插值曲线都行；边界条件（始末位形 + 速度）经 \cref{eq:nh-flat} 折算成平坦输出及其导数的端点值，用多项式/样条一插即得，再代数回算保证可行。这就是 \cref{sec:nh-steer} 多项式转向的全部底气。代价是平坦输出在奇点（独轮车 $\dot x=\dot y=0$，即停车—原地转）处退化，须特殊处理。
\end{insight}

% ════════════════════════════════════════════════════════════════
\section{开环转向：正弦输入与多项式轨迹}
\label{sec:nh-steer}

有了链式型与平坦，把系统从 $\mathbf{q}_i$ 驱到 $\mathbf{q}_f$ 的\emph{开环}转向就有了两类系统化方法。它们对应 \cref{thm:nh-three-ways} 的第三条出路（前馈），是 RRT、Hybrid A$^\star$ 等采样规划器所需的\emph{局部连接器}（steering function）的理论核心。

\subsection{正弦转向：Murray--Sastry 的优雅结果}
\label{sec:nh-steer-sin}

\paragraph{为何是正弦。}
考虑最小的非完整积分器（Brockett 系统 \cref{eq:nh-brockett-sys} 的两输入对称版）
\begin{equation}
  \dot q_1=u_1,\qquad \dot q_2=u_2,\qquad \dot q_3=q_1 u_2-q_2 u_1 .
  \label{eq:nh-integrator}
\end{equation}
要在 $t\in[0,1]$ 内从原点驱到 $(0,0,a)$（即只改变“面积型”状态 $q_3$、而 $q_1,q_2$ 复位），并最小化控制能量 $\int_0^1\lVert\mathbf{u}\rVert^2\mathrm{d}t$。这是个最优控制问题，Euler--Lagrange 方程（\cref{ch:optimal-control-basics}）给出最优输入满足 $\dot{\mathbf{u}}=\boldsymbol{\Lambda}\mathbf{u}$（$\boldsymbol{\Lambda}$ 反对称、为常数 Lagrange 乘子），解为 $\mathbf{u}(t)=e^{\boldsymbol{\Lambda}t}\mathbf{u}(0)$——即\emph{整数倍频率的正弦与余弦}。$q_3$ 的净变化恰是 $(u_1,u_2)$ 在相平面上扫过的有向面积，正弦输入把它做成一个椭圆环路，环路面积 $=a$。这便是“正弦转向”的来历\cite{murray1993sinusoids}。

\begin{theorem}[链式型的正弦/逐级转向\cite{murray1993sinusoids}]\label{thm:nh-sin-steer}
对 $(2,n)$ 链式型 \cref{eq:nh-chained}，可\emph{逐级}地把状态驱到目标：
\begin{enumerate}
  \item 用任意输入把直接受控的 $z_1,z_2$（零阶变量）驱到目标值；
  \item 再用一对整数倍频率的正弦 $v_1=a_1\sin(2\pi t),\ v_2=a_2\cos(2\pi t)$ 把一阶李括号方向的状态 $z_3$ 驱到目标，而由频率的整数关系，$z_1,z_2$ 在一个周期后\emph{自动复位}、不受影响；
  \item 更高阶变量 $z_4,\dots,z_n$ 用更高整数倍频率 $v_1=b\sin(2\pi k t)$ 类推，逐级钉死，低阶不被扰动。
\end{enumerate}
对一般 $m$ 输入的一阶系统（增长向量 $(m,m(m{+}1)/2)$，其代数来由见 \cref{sec:nh-lie-structure} 与 \cref{ins:nh-growth-firstorder}），当 $Y$ 满秩、$m$ 为偶数时，最优输入恰含 $m/2$ 个频率为 $2\pi,2\cdot2\pi,\dots,(m/2)\cdot2\pi$ 的正弦——Brockett 的“惊艳事实”。
\end{theorem}

\begin{insight}[逐级转向的节律：先钉零阶、再钉李括号]\label{ins:nh-sin-rhythm}
正弦转向的算法骨架是“\emph{按李括号阶数从低到高，一级一级钉死}”：直接可控的方向先随便走到位；要靠一阶李括号凑的方向，用一对整数倍频正弦走一个闭环把它挪到位、同时让已钉好的低阶变量绕回原值；二阶方向再用更高频正弦……频率的整数关系保证“高阶机动不破坏低阶成果”。这与 \cref{ins:nh-brockett-geo} 里“李括号方向只能靠二阶机动净挪”是同一回事的建设性版本——Brockett 说“瞬时使不上劲”，正弦转向说“那就用周期机动凑”。
\end{insight}

\subsection{多项式 / 平坦转向}
\label{sec:nh-steer-poly}

更直接的一路是借\emph{平坦}（\cref{sec:nh-flat}）：既然平坦输出空间无微分约束，就在那里插一条满足端点条件的多项式，再代数回算。

\begin{example}[独轮车的笛卡尔三次多项式转向]\label{ex:nh-poly-steer}
要把独轮车从 $\mathbf{q}_i=(x_i,y_i,\theta_i)$ 驱到 $\mathbf{q}_f=(x_f,y_f,\theta_f)$（参数 $s\in[0,1]$，非弧长）。平坦输出取 $(x,y)$，用三次多项式插值，并令端点切向匹配朝向：
\begin{equation}
  \begin{aligned}
    x(s)&=(1-s)^2(\,\cdots\,)\quad\text{满足 } x(0)=x_i,\ x(1)=x_f,\\
    x'(0)&=k_i\cos\theta_i,\quad x'(1)=k_f\cos\theta_f,
  \end{aligned}
  \qquad (y\text{ 同理，切向用 }\sin\theta)
  \label{eq:nh-poly}
\end{equation}
其中 $k_i,k_f>0$ 调形（影响出发/到达的“冲量”）。由 $(x(s),y(s))$ 经 \cref{eq:nh-flat-recon}（把 $\dot{}\to{}'$，再乘 $s$ 的时间律）得 $\theta(s)$ 与几何输入 $\tilde v,\tilde\omega$，自动满足非完整约束。端点切向取 $\theta_i,\theta_f$ 正是为了让朝向在两端对上。\emph{注意}：纯重定向（$x_i{=}x_f,y_i{=}y_f$ 但 $\theta_i\neq\theta_f$，即原地转）会触发平坦输出奇异，须改用经链式型坐标 \cref{eq:nh-uni2chain} 的参数化输入，或人为引入一个绕行 via 点。
\end{example}

\begin{table}[H]\centering\small
\caption{两类开环转向的对照}
\label{tab:nh-steer}
\begin{tabular}{lll}
\toprule
 & 正弦转向（\cref{sec:nh-steer-sin}） & 多项式/平坦转向（\cref{sec:nh-steer-poly}） \\
\midrule
工作空间 & 链式型坐标 $z$ & 平坦输出 $(x,y)$ 或链式型 $z_1,z_n$ \\
核心机理 & 整数倍频正弦逐级钉变量 & 平坦输出空间无约束插值 + 代数回算 \\
输出曲线 & 闭环式机动（含来回） & 直接的笛卡尔光滑路径 \\
奇异/退化 & 频率须整数关系 & 原地转（$\dot x{=}\dot y{=}0$）退化，需特判 \\
天然搭档 & 拖挂链等高非完整度系统 & 独轮/自行车的 RRT steering、轨迹生成 \\
\bottomrule
\end{tabular}
\end{table}

\begin{note}[还有第三类：参数化输入与时间最优]\label{note:nh-steer-more}
除正弦外，链式型在\emph{分段常值输入}下也可闭式积分（对应最优控制的 bang--bang 解），把规划化为求若干切换时刻；这是采样规划里 state-lattice 运动基元的来源。另一支是\emph{时间最优}轨迹：对“前进速度有界、可前可后”的小车，最短时间路径由有限种基元弧（直线段 + 最小半径圆弧，段间可含尖点/换向）拼成（Reeds--Shepp 曲线即属此族），段数与弧长有结构性上界。这些属规划部 P3 的范畴，本章只点到接口。
\end{note}

% ════════════════════════════════════════════════════════════════
\section{轨迹跟踪：可行轨迹是可镇定的}
\label{sec:nh-track}

现在转入闭环。\textbf{轨迹跟踪}要机器人渐近跟上一条\emph{随时间运动}的参考轨迹 $(x_d(t),y_d(t))$。关键前提：该参考轨迹必须对运动学模型\emph{可行}（admissible）——存在参考输入 $v_d,\omega_d$ 使
\begin{equation}
  \dot x_d=v_d\cos\theta_d,\quad \dot y_d=v_d\sin\theta_d,\quad \dot\theta_d=\omega_d,
  \qquad \theta_d=\operatorname{Atan2}(\dot y_d,\dot x_d),
  \label{eq:nh-ref}
\end{equation}
即 $\theta_d$ 由平坦性（\cref{ex:nh-flat-unicycle}）从笛卡尔轨迹算出。只要轨迹来自 \cref{sec:nh-steer} 的规划，便天然可行。与定点镇定不同，\textbf{跟踪不撞 Brockett 障碍}（参考持续运动，见前文 pitfall），故有干净的光滑（时变）反馈解。

\subsection{机器人系下的跟踪误差}
\label{sec:nh-track-err}

直接用 $\mathbf{q}_d-\mathbf{q}$ 作误差不好（笛卡尔误差与朝向误差耦合）。标准做法是把误差转到\emph{机器人体坐标系}：
\begin{equation}
  \mathbf{e}=\begin{bmatrix}e_1\\e_2\\e_3\end{bmatrix}
  =\begin{bmatrix}\cos\theta&\sin\theta&0\\-\sin\theta&\cos\theta&0\\0&0&1\end{bmatrix}
  \begin{bmatrix}x_d-x\\y_d-y\\\theta_d-\theta\end{bmatrix}.
  \label{eq:nh-track-err}
\end{equation}
$e_1$ 是纵向（前后）误差、$e_2$ 是横向（左右）误差、$e_3$ 是朝向误差。对 \cref{eq:nh-track-err} 求导并代入独轮车方程，得误差动力学
\begin{equation}
  \begin{aligned}
    \dot e_1&=\omega e_2-v+v_d\cos e_3,\\
    \dot e_2&=-\omega e_1+v_d\sin e_3,\\
    \dot e_3&=\omega_d-\omega .
  \end{aligned}
  \label{eq:nh-track-dyn}
\end{equation}
为分离前馈与反馈，引入输入偏差 $u_1=v_d\cos e_3-v,\ u_2=\omega_d-\omega$，则第一、三式线性化，第二式仍非线性且时变（含 $v_d(t),\omega_d(t)$）。

\subsection{近似线性化控制器}
\label{sec:nh-track-lin}

最简单的设计：在参考轨迹（$\mathbf{e}=\mathbf{0}$）附近线性化误差动力学——令 $\sin e_3\approx e_3$、矩阵在轨迹上取值，得
\begin{equation}
  \dot{\mathbf{e}}=\begin{bmatrix}0&\omega_d&0\\-\omega_d&0&v_d\\0&0&0\end{bmatrix}\mathbf{e}
  +\begin{bmatrix}1&0\\0&0\\0&1\end{bmatrix}\begin{bmatrix}u_1\\u_2\end{bmatrix}.
  \label{eq:nh-track-linsys}
\end{equation}
这是\emph{时变}线性系统（$v_d,\omega_d$ 随时间）。取线性反馈
\begin{equation}
  u_1=-k_1 e_1,\qquad u_2=-k_2 e_2-k_3 e_3,
  \label{eq:nh-track-linfb}
\end{equation}
其中增益按“把闭环特征多项式配成一个二阶 $+$ 一阶”的标准形选取——令期望极点为 $-2\zeta a$ 与一对阻尼 $\zeta$、自然频率 $a$ 的共轭复根，得
\begin{equation}
  k_1=k_3=2\zeta a,\qquad k_2=\frac{a^2-\omega_d^2}{v_d}\quad(\text{当 } v_d\neq0).
  \label{eq:nh-track-gains}
\end{equation}
$k_2$ 含 $v_d$，故控制律本身\emph{时变}。当 $\mathbf{e}\to\mathbf{0}$ 时实际输入 $v\to v_d,\omega\to\omega_d$，反馈退化为纯前馈。

\begin{pitfall}[时变系统：极点负实部 $\neq$ 渐近稳定]
\cref{eq:nh-track-linsys} 是时变系统，把它“冻结”后算出的特征值即便处处有负实部，也\emph{不能}直接断言渐近稳定——时变线性系统的稳定性不由瞬时特征值决定。严格的稳定性须在 $v_d,\omega_d$ 不同时趋零（\emph{持续激励}）等假设下另证。这正是“跟踪可、镇定难”的技术映照：一旦参考停下（$v_d,\omega_d\to0$），\cref{eq:nh-track-gains} 的 $k_2$ 失定义、\cref{eq:nh-track-linsys} 退化为不可控，跟踪框架自然失效——回到 Brockett 障碍。
\end{pitfall}

\subsection{非线性（Lyapunov）控制器}
\label{sec:nh-track-nl}

要全局收敛、且不依赖小误差近似，回到精确误差动力学 \cref{eq:nh-track-dyn}，取控制律
\begin{equation}
  u_1=-k_1(v_d,\omega_d)\,e_1,\qquad u_2=-k_2\,v_d\,\frac{\sin e_3}{e_3}\,e_2-k_3(v_d,\omega_d)\,e_3,
  \label{eq:nh-track-nlfb}
\end{equation}
其中 $k_1,k_3>0$ 为有界且导数有界的函数、$k_2>0$ 常数。

\begin{theorem}[非线性跟踪的全局收敛\cite{siciliano2009robotics}]\label{thm:nh-track-nl}
设参考输入 $v_d,\omega_d$ 有界、导数有界，且\emph{不同时}收敛到零（持续激励）。则 \cref{eq:nh-track-nlfb} 使跟踪误差 $\mathbf{e}$ 全局收敛到零（任意初值）。
\end{theorem}

\begin{proof}[Lyapunov + Barbalat 思路]
取 $V=\tfrac12(e_1^2+e_2^2)+\tfrac{1}{2k_2}e_3^2\ge0$。沿闭环求导，交叉项 $\omega e_1 e_2$ 抵消、前馈项经控制律配平（控制律 \cref{eq:nh-track-nlfb} 中的 $\sin e_3/e_3$ 因子恰把 $\dot e_2$ 里的 $v_d\sin e_3$ 与 $\tfrac1{k_2}e_3\dot e_3$ 配成 $-\tfrac{k_3}{k_2}e_3^2$），得 $\dot V=-k_1 e_1^2-\tfrac{k_3}{k_2}e_3^2\le0$（负半定）。故 $V$ 有下界且单调不增、收敛到极限，$\lVert\mathbf{e}\rVert$ 有界。因系统时变，\emph{不能}用 LaSalle；改用 \textbf{Barbalat 引理}：验证 $\ddot V$ 有界 $\Rightarrow\dot V$ 一致连续 $\Rightarrow\dot V\to0$，于是 $e_1,e_3\to0$。再回代系统方程并用持续激励（$v_d^2+\omega_d^2\not\to0$），推出 $e_2\to0$。$\square$
\end{proof}

\begin{note}[Barbalat 引理：时变系统的“准 LaSalle”]\label{note:nh-barbalat}
\textbf{Barbalat 引理}：若 $\phi(t)$ 一致连续且 $\int_0^\infty\phi\,\mathrm{d}t$ 存在有限，则 $\phi(t)\to0$。对 Lyapunov 分析，常用其推论：$V$ 下有界、$\dot V\le0$、$\ddot V$ 有界（$\Rightarrow\dot V$ 一致连续），则 $\dot V\to0$。它是\emph{时变}非线性系统稳定性分析的主力工具——LaSalle 不变集原理只对\emph{自治}系统成立，而跟踪误差动力学因含 $v_d(t),\omega_d(t)$ 而时变，故本章跟踪与镇定的收敛证明几乎都落到 Barbalat 上（\cref{sec:ctrl-lyapunov} 给出基础）。
\end{note}

\subsection{输入/输出线性化：跟踪车前一点}
\label{sec:nh-track-io}

若不要求跟踪\emph{接触点}本身、而允许跟踪车体上一个偏置点，可得彻底解耦的线性控制。取输出为纵轴上距接触点 $|b|$ 处的点 $B$（$b\neq0$）：
\begin{equation}
  y_1=x+b\cos\theta,\qquad y_2=y+b\sin\theta .
  \label{eq:nh-io-out}
\end{equation}
求导得 $\dot{\mathbf{y}}=\mathbf{T}(\theta)(v,\omega)^\top$，$\mathbf{T}=\begin{bsmallmatrix}\cos\theta&-b\sin\theta\\\sin\theta&b\cos\theta\end{bsmallmatrix}$，$\det\mathbf{T}=b\neq0$ 处处可逆。用输入变换 $(v,\omega)^\top=\mathbf{T}^{-1}(\theta)(u_1,u_2)^\top$ 把系统化为
\begin{equation}
  \dot y_1=u_1,\qquad \dot y_2=u_2,\qquad \dot\theta=\frac{u_2\cos\theta-u_1\sin\theta}{b},
  \label{eq:nh-io-lin}
\end{equation}
即两个解耦积分器 $+$ 一段不受控的朝向内动态。简单线性律 $u_i=\dot y_{id}+k_i(y_{id}-y_i)$（$k_i>0$）即得点 $B$ 笛卡尔误差的\emph{指数解耦}收敛。

\begin{insight}[偏置点的代价：朝向不受控 + 扫掠区域]\label{ins:nh-io-tradeoff}
I/O 线性化把非完整难点“藏进”了不受控的朝向内动态（\cref{eq:nh-io-lin} 第三式）。好处：参考路径可\emph{任意}，甚至带折角（如方形）也不必停车重定向——因为受控的是\emph{车前 $b$ 处}的点 $B$，不是接触点。代价有二：其一，车体朝向不被控制；其二，$B$ 走折角时\emph{接触点}走一条“抹圆”的内侧曲线，车身会扫掠出额外区域，须留无碰撞通道。$b\to0$ 时跟踪越准、但 $\mathbf{T}$ 趋奇异，转向速度在折角处暴涨、易饱和。$b$ 是“跟踪精度 vs 控制平顺”的旋钮。
\end{insight}

\subsection{动态反馈线性化}
\label{sec:nh-track-dfl}

更彻底的精确跟踪用\emph{动态}反馈线性化：把前向速度 $v$ 升格为状态、引入其导数为新输入（加一阶积分器 $\dot v=\xi$），可使\emph{接触点本身}的 $(x,y)$ 动态精确线性化为双积分器，从而精确跟踪任意 $C^2$ 笛卡尔轨迹（含朝向一致）。其代价是控制律含 $1/\xi$（$\xi$ 为升格后的速度状态），当 $\xi\to0$ 时\emph{转向输入 $\omega$ 无界}。

\begin{pitfall}[动态反馈线性化不能用于定点镇定]
$\omega\propto1/\xi$ 意味着\emph{机器人一旦停下（$\xi\to0$）控制律就奇异}。因此动态反馈线性化要求“永不停车”，\emph{无法}用于把车停到一个固定姿态——它又一次撞上 Brockett 障碍。这与 \cref{thm:nh-track-nl}（需持续激励）、\cref{eq:nh-track-gains}（$k_2$ 含 $1/v_d$）一脉相承：本节四个跟踪控制器\emph{无一}能兼做定点镇定。定点镇定必须另起炉灶（\cref{sec:nh-reg}）。
\end{pitfall}

\begin{table}[H]\centering\small
\caption{四种跟踪控制器对照（皆要求参考轨迹可行且持续激励）}
\label{tab:nh-track}
\begin{tabular}{llll}
\toprule
方法 & 受控对象 & 收敛性 & 停车（$v_d\to0$）时 \\
\midrule
近似线性化（\cref{sec:nh-track-lin}） & 全状态 $\mathbf{e}$ & 局部、时变 & $k_2$ 失定义 \\
非线性 Lyapunov（\cref{sec:nh-track-nl}） & 全状态 $\mathbf{e}$ & 全局（持续激励） & 失效（需 $v_d^2{+}\omega_d^2\not\to0$） \\
I/O 线性化（\cref{sec:nh-track-io}） & 偏置点 $B$，朝向不控 & 指数、解耦 & 仍可（但不控朝向） \\
动态反馈线性化（\cref{sec:nh-track-dfl}） & 接触点 $(x,y)$ 精确 & 指数 & 奇异（$\omega\to\infty$） \\
\bottomrule
\end{tabular}
\end{table}

% ════════════════════════════════════════════════════════════════
\section{定点姿态镇定：Brockett 障碍的正面交锋}
\label{sec:nh-reg}

最后处理\textbf{姿态镇定}（posture regulation）：把机器人从任意初始位形驱到一个\emph{固定}目标位形 $\mathbf{q}_d$ 并停住（不设定逼近路径）。由 \cref{sec:nh-brockett}，这是非完整控制最难的一关——\cref{sec:nh-track} 的跟踪控制器\emph{无一}可用（皆需持续激励），且不存在光滑时不变反馈（Brockett）。这里没有“通用控制器”，须按 \cref{thm:nh-three-ways} 设计特设律。不失一般性设 $\mathbf{q}_d=\mathbf{0}$。

\begin{theorem}[非完整系统无通用镇定器]\label{thm:nh-no-universal}
独轮车（及一般非完整系统）\emph{不}存在“通用控制器”——即不存在能渐近镇定任意（持续或非持续）状态轨迹的单一反馈律。这与机械臂形成鲜明对比：逆动力学控制对机械臂就是通用镇定器。根源是非完整性：定点（非持续）目标违反 Brockett 必要条件（\cref{thm:nh-brockett}）。故姿态镇定必须用\emph{专门设计}的、时变或不连续的控制律。
\end{theorem}

\subsection{笛卡尔镇定：只管到点，不管朝向}
\label{sec:nh-reg-cart}

先解一个更易、且实用的退化版：只把机器人驱到目标\emph{位置}（不指定终端朝向）。例如装环形传感器的机器人巡视一串视点，到点即可、朝向无所谓。设目标在原点、笛卡尔误差 $\mathbf{e}_p=(-x,-y)$，取
\begin{equation}
  v=-k_1(x\cos\theta+y\sin\theta),\qquad \omega=k_2\big(\operatorname{Atan2}(y,x)-\theta+\pi\big),
  \label{eq:nh-reg-cart}
\end{equation}
$k_1,k_2>0$。几何意义直白：$v$ 正比于位置误差在车体纵轴上的投影（误差在前就前进、在后就倒车），$\omega$ 正比于车头朝向与“指向误差矢量”之间的夹角（\emph{指向误差}），把车头转向目标。

\begin{proof}[半正定 Lyapunov + Barbalat]
取 $V=\tfrac12(x^2+y^2)$（在原点仅\emph{半正定}：$x{=}y{=}0$ 时无论 $\theta$ 取何值 $V$ 皆零）。代 \cref{eq:nh-reg-cart} 与独轮车方程得 $\dot V=-k_1(x\cos\theta+y\sin\theta)^2\le0$。$V$ 下有界、$\ddot V$ 有界 $\Rightarrow$ Barbalat $\Rightarrow\dot V\to0$，即位置误差在纵轴上的投影 $x\cos\theta+y\sin\theta\to0$。这只能发生在原点：否则 $\omega$ 律会持续把车头转向误差矢量、使投影非零。故 $\mathbf{e}_p\to0$，任意初值。$\square$
\end{proof}

\subsection{全姿态镇定：极坐标与 $\sigma$ 过程}
\label{sec:nh-reg-polar}

要连\emph{朝向}也镇定，关键技巧是改用\emph{极坐标}描述误差。设目标为原点，定义
\begin{equation}
  \rho=\sqrt{x^2+y^2},\qquad
  \gamma=\operatorname{Atan2}(y,x)-\theta+\pi,\qquad
  \delta=\gamma+\theta ,
  \label{eq:nh-polar}
\end{equation}
$\rho$ 为到原点距离、$\gamma$ 为“误差矢量与车体纵轴的夹角”、$\delta$ 为“误差矢量与 $x$ 轴的夹角”。独轮车在极坐标下成为
\begin{equation}
  \dot\rho=-v\cos\gamma,\qquad
  \dot\gamma=\frac{\sin\gamma}{\rho}\,v-\omega,\qquad
  \dot\delta=\frac{\sin\gamma}{\rho}\,v .
  \label{eq:nh-polar-dyn}
\end{equation}

\begin{theorem}[极坐标姿态镇定律\cite{siciliano2009robotics}]\label{thm:nh-reg-polar}
取反馈
\begin{equation}
  v=k_1\rho\cos\gamma,\qquad
  \omega=k_2\gamma+k_1\frac{\sin\gamma\cos\gamma}{\gamma}\,(\gamma+k_3\delta),
  \label{eq:nh-reg-polar-fb}
\end{equation}
$k_1,k_2>0$。则极坐标系统 \cref{eq:nh-polar-dyn} 渐近收敛到 $(\rho,\gamma,\delta)=(0,0,0)$，即机器人到达目标位形（位置 $+$ 朝向）。
\end{theorem}

\begin{proof}[思路]
取 $V=\tfrac12(\rho^2+\gamma^2+k_3\delta^2)$，沿闭环得 $\dot V=-k_1\cos^2\gamma\,\rho^2-k_2\gamma^2\le0$（负半定）。$V$ 收敛、状态有界；$\ddot V$ 有界 $\Rightarrow$ Barbalat $\Rightarrow\dot V\to0\Rightarrow\rho,\gamma\to0$；再分析闭环知 $\delta\to0$。注意 \cref{eq:nh-reg-polar-fb} 用了 $\sin\gamma\cos\gamma/\gamma$，$\gamma\to0$ 时极限为 $\cos\gamma\to1$，连续可延拓。$\square$
\end{proof}

\begin{insight}[$\sigma$ 过程：把奇点“吹散”再镇定]\label{ins:nh-sigma}
\cref{eq:nh-polar} 这步“笛卡尔 $\to$ 极坐标”正是非完整镇定经典的 $\sigma$ \emph{过程}（坐标奇异变换 / blow-up）：在原点处把一个点“吹胀”成一族方向，使原本无法被光滑反馈处理的退化点，在新坐标里变得可控可镇。角 $\gamma,\delta$ 在 $x{=}y{=}0$ 处无定义（原点被吹散），但在\emph{运动过程中}始终良定，并渐近趋于目标零值。代价正是 Brockett 障碍的“收据”：\cref{eq:nh-reg-polar-fb} 映回原始笛卡尔坐标后，在原点处\emph{不连续}——这与“任何能镇定独轮车姿态的反馈\emph{必}对状态不连续和/或时变”完全吻合。$\sigma$ 过程是把这种不可避免的不连续性\emph{安置}在一个无害位置（原点）的优雅手法。
\end{insight}

\begin{example}[切换控制：另一条不连续出路]\label{ex:nh-switch}
对独轮车，一个更朴素的不连续方案是\emph{分阶段切换}：先原地转使车头指向目标（只动 $\omega$）；再直行到目标位置（只动 $v$）；最后原地转到目标朝向（只动 $\omega$）。三段每段都是平凡可镇的，拼起来把任意初位形驱到 $\mathbf{q}_d$。它显式地放弃了光滑性（段间切换不连续），同样绕过 Brockett 障碍（\cref{thm:nh-three-ways} 第二条）。工程上它鲁棒、易实现，缺点是含原地转、不平顺，且对 $\rho\to0$ 附近的朝向抖动敏感。
\end{example}

\begin{table}[H]\centering\small
\caption{跟踪 vs 镇定：非完整控制的两类任务总览}
\label{tab:nh-track-vs-reg}
\begin{tabular}{lll}
\toprule
 & 轨迹跟踪（\cref{sec:nh-track}） & 定点姿态镇定（\cref{sec:nh-reg}） \\
\midrule
目标 & 跟上运动参考 $(x_d(t),y_d(t))$ & 停到固定位形 $\mathbf{q}_d$ \\
撞 Brockett 障碍？ & 否（参考持续运动） & \textbf{是}（目标静止） \\
光滑时不变反馈？ & 可（时变线性即够） & \textbf{不可}（\cref{thm:nh-brockett}） \\
典型方法 & 近似/非线性/IO/动态 FBL & 时变 / 不连续（极坐标 $\sigma$、切换） \\
工程地位 & 主力（避障靠预规划轨迹） & 末端泊车 / 对接 \\
\bottomrule
\end{tabular}
\end{table}

\begin{practice}
\begin{enumerate}
  \item 验证 \cref{eq:nh-uni2chain} 的输入变换可逆，并写出由 $(v_1,v_2)$ 反解 $(v,\omega)$ 的表达式；指出在何处奇异。
  \item 对 Brockett 系统 \cref{eq:nh-brockett-sys}，自行补全 \cref{ex:nh-brockett-fail} 的论证：证明对任意 $a\neq0$，$\mathbf{f}(\mathbf{x},\mathbf{u})=(0,0,a)^\top$ 无解，故不满足 Brockett 必要条件。
  \item 用平坦输出 $(x,y)$ 为独轮车规划一条从 $(0,0,0)$ 到 $(2,1,\pi/2)$ 的三次多项式轨迹（\cref{ex:nh-poly-steer}），写出 $\theta(s),v(s),\omega(s)$ 的表达式，并检查端点朝向是否对上。
  \item 解释为什么 \cref{sec:nh-track} 的四种跟踪控制器都\emph{不能}用于把独轮车停到一个固定姿态；分别指出每个控制器在 $v_d\to0$（或 $\xi\to0$）时具体出什么问题。
  \item （开放题）极坐标镇定律 \cref{eq:nh-reg-polar-fb} 映回笛卡尔坐标后在原点不连续。试论证：这种不连续性是\emph{不可避免}的（提示：Brockett \cref{thm:nh-brockett} $+$ \cref{thm:nh-no-universal}）。
\end{enumerate}
\end{practice}

% ════════════════════════════════════════════════════════════════
\section*{本章小结}
\addcontentsline{toc}{section}{本章小结}

本章把非完整系统的\emph{运动规划与控制}半边讲透，与 \cref{ch:mm-wheeled} 的\emph{建模与可控性}半边互补。主线一句话：\textbf{非完整系统可控（Chow）却不可光滑镇定（Brockett），于是规划靠链式型 + 平坦，控制分跟踪（易）与定点镇定（难）两途。}

\begin{enumerate}
  \item \textbf{定位}（\cref{sec:nh-place}）：非完整 = 速度约束不可积，唯一“瞬时受限、整体可达”者；与欠驱动（执行器少）、完整（位形受困）三足鼎立。
  \item \textbf{Brockett 障碍}（\cref{sec:nh-brockett}）：光滑时不变状态反馈要求 $\mathbf{f}$ 的像含原点邻域；非完整系统在李括号方向上瞬时速度为零，过不了这关，故\emph{无法被光滑反馈镇定到一点}，只剩时变 / 不连续 / 开环三条路。
  \item \textbf{链式型 + 平坦}（\cref{sec:nh-chained,sec:nh-flat}）：$(2,n)$ 链式型是两输入非完整系统的规范形，可控且非完整度 $\kappa=n-1$；无漂移系统“平坦 $\Leftrightarrow$ 可化链式型”，独轮车笛卡尔坐标即平坦输出，可行轨迹\emph{代数地}回算、近乎免费。
  \item \textbf{开环转向}（\cref{sec:nh-steer}）：正弦转向（整数倍频、逐级钉变量，Murray--Sastry）与多项式/平坦转向（输出空间无约束插值），是采样规划器的局部连接器。
  \item \textbf{轨迹跟踪}（\cref{sec:nh-track}）：参考持续运动则不撞 Brockett；机器人系误差 $+$ 近似线性化 / 非线性 Lyapunov（Barbalat）/ I/O 线性化（跟车前一点）/ 动态反馈线性化——皆可镇定，但无一能兼做定点镇定。
  \item \textbf{定点镇定}（\cref{sec:nh-reg}）：无通用镇定器；笛卡尔镇定（半正定 Lyapunov）只管到点，全姿态镇定靠极坐标 $\sigma$ 过程（映回笛卡尔必在原点不连续）或分段切换——是 Brockett 障碍的正面体现。
\end{enumerate}

\noindent\textbf{承上启下}：本章上承欠驱动（\cref{ch:underactuated}，同为受限但机理不同）与线性可控性（\cref{sec:ctrl-controllability}，本章是其非线性几何推广）；李括号、可控性的几何根基见 \cref{ch:lie,ch:geometric-mechanics}，底盘建模与 Chow/LARC 细节见 \cref{ch:mm-wheeled}。下接：开环转向是规划部 P3 采样/格点规划器的 steering 内核；轮式底盘的非完整约束在轮足/移动操作（P7、P9）里与浮基动力学、全身控制汇流，届时“瞬时不能横移”将与“足端接触约束”在统一 Pfaffian 框架下并肩处理。

\section*{延伸阅读}
\addcontentsline{toc}{section}{延伸阅读}
\begin{itemize}
  \item 几何控制理论与正弦转向的系统阐述：Murray, Li, Sastry《A Mathematical Introduction to Robotic Manipulation》第 7--8 章\cite{murray1994mathematical}；原始正弦转向论文 Murray \& Sastry (1993)\cite{murray1993sinusoids}。
  \item 轮式机器人建模、链式型、平坦与跟踪/镇定的工程化处理：Siciliano 等《Robotics: Modelling, Planning and Control》第 11 章\cite{siciliano2009robotics}；Spong 等《Robot Modeling and Control》第 14 章\cite{spong2006robot}（Brockett 障碍与动态反馈线性化的清晰陈述）。
  \item Brockett 必要条件原文：Brockett (1983)《Asymptotic stability and feedback stabilization》\cite{brockett1983asymptotic}。
  \item 采样规划与非完整 steering 的工程接口：LaValle《Planning Algorithms》\cite{lavalle2006planning}。
\end{itemize}

% ════════════════════════════════════════════════════════════════
\section{最优路径：Dubins 车与 Reeds--Shepp 曲线}
\label{sec:nh-optimal-paths}

\cref{sec:nh-steer} 的正弦与多项式转向解决了“怎么把系统从 $\mathbf{q}_i$ 驱到 $\mathbf{q}_f$”，但没问“最短怎么走”。本节给\cref{thm:nh-three-ways} 第三条出路（开环前馈）配上\emph{最优性}：在“前向速度恒定、转弯半径有下界”的几何约束下，$\SE(2)$ 上两位形之间的\emph{最短}可行路径有惊人简单的结构——只由\emph{最大曲率圆弧}与\emph{直线段}拼成，且候选词型有限。这就是 Dubins 车（只前进）与 Reeds--Shepp 车（容许倒车）两族曲线。它们是 \cref{note:nh-steer-more} 末尾点到的“时间最优基元弧”的展开，也是采样规划器（RRT、Hybrid A$^\star$）局部连接函数与非对称距离度量的解析内核。

\begin{remark}[与已有处理的分工]\label{rem:nh-optpath-scope}
这两族曲线在本书出现三次，各取一角，互不重复：规划部以它们作 \texttt{Steer} 与距离度量（六词型陈述、RRT 集成、非对称性，\cref{sec:plan-dubins}）；最优控制部由 Pontryagin 极小值原理\emph{推出}词型的有限性（共态结构、bang--bang，\cref{thm:pmp-dubins}）；\textbf{本节}站在\emph{非完整几何}的立场，讲它们\emph{为何是这条曲率受限约束下的最短路}、Reeds--Shepp\emph{充分族的精确结构}，以及它们与本章\emph{微分平坦 / 弧长参数化}的接口。三处宜对照阅读。
\end{remark}

\subsection{Dubins 车：只前进的曲率受限最短路}
\label{sec:nh-dubins}

\paragraph{模型。}
\textbf{Dubins 车}\cite{dubins1957curves}是恒定前向速率、转弯半径有下界 $\rho$（即曲率上界 $1/\rho$）的独轮车——正是 \cref{eq:nh-bicycle} 自行车在转向打死 $\phi=\pm\phi_{\max}$ 时给出的最小半径约束的理想化：
\begin{equation}
  \dot x=v\cos\theta,\quad \dot y=v\sin\theta,\quad \dot\theta=\frac{v}{\rho}\,u,\quad u\in[-1,1],\ v\equiv\text{const}>0 .
  \label{eq:nh-dubins-model}
\end{equation}
$u=+1/-1$ 对应以最小半径 $\rho$ 左/右转，$u=0$ 对应直行。求\emph{带朝向}的两位形 $\mathbf{q}_i,\mathbf{q}_f\in\SE(2)$ 间的最短可行路径。

\begin{theorem}[Dubins 最短路：六词型]\label{thm:nh-dubins}
曲率上界 $1/\rho$ 约束下，$\SE(2)$ 上任意两位形间的最短前进路径必由最大曲率圆弧（左转 $\mathrm{L}$、右转 $\mathrm{R}$，半径恒为 $\rho$）与直线段（$\mathrm{S}$）拼成，且属于下列六词型之一：
\begin{equation}
  \underbrace{\{\,\mathrm{LSL},\ \mathrm{LSR},\ \mathrm{RSL},\ \mathrm{RSR}\,\}}_{\text{CSC 型（含直线段）}}\ \cup\
  \underbrace{\{\,\mathrm{LRL},\ \mathrm{RLR}\,\}}_{\text{CCC 型（三弧相切）}},
  \label{eq:nh-dubins-words}
\end{equation}
其中 C 为最大曲率弧、S 为直线段。求最短路只需对六词型各算一次闭式长度、取最小。
\end{theorem}

\begin{insight}[词型有限性的根：bang--bang $+$ 共态]\label{ins:nh-dubins-why}
\cref{eq:nh-dubins-model} 中归一化转向 $u$\emph{线性}进入哈密顿量，故时间最短的最优控制是 bang--bang：$u^\star=\operatorname{sign}(\lambda_\theta)=\pm1$（最大曲率弧）或在奇异弧 $\lambda_\theta\equiv0$ 上 $u=0$（直线）。位置共态 $\lambda_x,\lambda_y$ 守恒、$\lambda_\theta$ 沿轨迹为正弦型，其零点至多让控制切换有限次，穷举即得六词型。这一推导即 Pontryagin 极小值原理在该系统上的范例，详见 \cref{thm:pmp-dubins}；本章只取其\emph{几何收成}——最短路是“弧—线—弧”的有限拼接，规划由无穷维降为六选一的代数比较。
\end{insight}

\paragraph{几何构造与闭式长度。}
六词型对应清晰的几何图像：在 $\mathbf{q}_i$ 处按 $\mathrm{L}/\mathrm{R}$ 各画一个半径 $\rho$ 的相切圆（“起圆”），在 $\mathbf{q}_f$ 处同样画两个“终圆”；\emph{CSC} 型取起圆与终圆的一条\emph{公切线}（外公切线对应同向 LSL/RSR，内公切线对应反向 LSR/RSL）作中间直线段，长度 $=$ 起弧长 $+$ 切线长 $+$ 终弧长，全部由圆心距与切点角\emph{闭式}给出；\emph{CCC} 型在两圆距离足够近时插入一个与两者皆相切的\emph{中间桥接圆}（反向转），适用于近距离大转角。各词型用统一的“以 $\rho$ 归一化、把 $\mathbf{q}_f$ 变到 $\mathbf{q}_i$ 体坐标系”后，长度只是起末相对位姿 $(\Delta x/\rho,\Delta y/\rho,\Delta\theta)$ 的初等三角函数，$O(1)$ 求值。

\begin{pitfall}[Dubins 距离非对称、且非度量]\label{pit:nh-dubins-metric}
因朝向\emph{有向}，Dubins 路径长 $d(\mathbf{q}_a,\mathbf{q}_b)\neq d(\mathbf{q}_b,\mathbf{q}_a)$（掉头要绕大圈）：它\emph{不}是对称度量。用作采样规划的最近邻查询时必须用这个非对称代价，不能退回欧氏 $\lVert\cdot\rVert$，否则 \texttt{Steer} 出的边物理不可行（\cref{sec:plan-dubins} 详述工程后果）。此外 Dubins 车不能倒车，故当 $\mathbf{q}_f$ 在 $\mathbf{q}_i$ 正后方近处时最短路可能很长（要兜一个大环）——这正是下面 Reeds--Shepp 容许倒车的动机。
\end{pitfall}

\subsection{Reeds--Shepp 曲线：容许倒车的完整族}
\label{sec:nh-reeds-shepp}

允许\emph{前进与倒车}（$v$ 可正可负、可在路径中途变号），即得\textbf{Reeds--Shepp 车}\cite{reeds1990optimal}，更贴合能倒车的汽车。代价是最优路径在方向反转处出现\emph{尖点}（cusp）——速度过零、车体瞬时停住再反向，而词型大幅增多。

\begin{definition}[尖点（cusp）：运动反向]\label{def:nh-cusp}
\emph{尖点}是 Reeds--Shepp 路径上前向速率 $v$ 变号（$v=0$）的孤立点，几何上对应笛卡尔路径 $(x(s),y(s))$ 的一个“折返”：曲线在该点处切向反转 $180^\circ$、形成一个尖。在尖点 $\bar s$ 处 $x'(\bar s)=y'(\bar s)=0$，由 \cref{ex:nh-flat-unicycle} 的平坦重构 \cref{eq:nh-flat-recon}，朝向公式 $\theta=\operatorname{Atan2}(\dot y,\dot x)$ 在此\emph{失定义}（$0/0$），须按\emph{单侧极限}$\theta(\bar s)=\lim_{s\to\bar s^-}\operatorname{Atan2}(y',x')$\emph{由连续性}补出（转向速率 $\omega$ 同理）。这与 \cref{ins:nh-flat-why} 末尾警告的“平坦输出奇点（停车—换向处）”是\emph{同一回事}：倒车换向恰是平坦坐标的退化点。
\end{definition}

\begin{theorem}[Reeds--Shepp 充分族：九组与 48 词型\cite{reeds1990optimal}]\label{thm:nh-rs}
设前进/倒车速度与转向速度有界、最小转弯半径恒为 $\rho=v_{\max}/\omega_{\max}$。则容许倒车的两位形间\emph{时间（亦即弧长）最短}路径，落在一个由\emph{两类基元弧}——半径 $\rho$ 的圆弧 $C$（可正/负向、可左/右转）与直线段 $S$（可正/负向）——拼接而成的\textbf{充分族}里。以 $C_a,S_e$ 记时长（$v_{\max}=\omega_{\max}=1$ 时即弧长）、上标 $\pm$ 记前进/倒车、$\{\zeta\}^\flat$ 记一处尖点，充分族可归为\textbf{九组}：
\begin{equation}
  \begin{array}{lll}
  \mathrm{I} & C_a\,|\,C_b\,|\,C_e & a,b,e\ge0,\ a+b+e\le\pi\\
  \mathrm{II},\mathrm{III} & C_a\,|\,C_b\,C_e,\ \ C_a\,C_b\,|\,C_e & 0\le a,e\le b\le\pi/2\\
  \mathrm{IV},\mathrm{V} & C_a\,C_b\,|\,C_b\,C_e,\ \ C_a\,|\,C_b\,C_b\,|\,C_e & 0\le a,e\le b\le\pi/2\\
  \mathrm{VI} & C_a\,|\,C_{\pi/2}\,S_e\,C_{\pi/2}\,|\,C_b & 0\le a,b\le\pi/2,\ e\ge0\\
  \mathrm{VII},\mathrm{VIII} & C_a\,|\,C_{\pi/2}\,S_e\,C_b,\ \ C_a\,S_e\,C_{\pi/2}\,|\,C_b & \\
  \mathrm{IX} & C_a\,S_e\,C_b & 0\le a,b\le\pi/2,\ e\ge0
  \end{array}
  \label{eq:nh-rs-groups}
\end{equation}
每条路径不超过\emph{五段}基元弧；圆弧时长被 $\pi/2$ 或 $\pi$ 界住，直线段长由 $\mathbf{q}_i,\mathbf{q}_f$ 的笛卡尔距离定。九组按运动方向与转向左右实例化，共生成 \textbf{48} 种序列，穷举各算闭式长度取最小即得最短路。
\end{theorem}

\begin{remark}[48 与 46：路径数的精确归属]\label{rem:nh-rs-count}
\cref{eq:nh-rs-groups} 的 48 是 Reeds--Shepp（1990）原文的\emph{充分}集；Sussmann--Tang（1991）用李代数方法把它缩到 \textbf{46} 条。常见误述把“46”错挂在 Reeds--Shepp 名下——精确表述见 \cref{thm:pmp-dubins} 下的辨析（“48 出自 1990、46 出自 1991”），此处不重复。工程上 Reeds--Shepp 距离是带倒车泊车、Hybrid A$^\star$ 等自动驾驶规划器的标准度量；固定翼无人机不能倒飞，则退回 Dubins。
\end{remark}

\begin{figure}[H]
\centering
\begin{tikzpicture}[>=Stealth, scale=0.95]
  % --- Dubins LSL (top): arc-line-arc, all forward, no cusp ---
  \coordinate (s) at (0,0);
  \draw[->, main!70!black, thick] (s) -- ++(0.6,0);
  \fill[main!70!black] (s) circle (1.3pt) node[below left, font=\scriptsize]{$\mathbf{q}_i$};
  \draw[second!45, dashed] (0,0.85) circle (0.85);
  \draw[main!70!black, very thick] (0,0) arc (-90:-22:0.85);
  \coordinate (t1) at ($(0,0.85)+(-22:0.85)$);
  \coordinate (t2) at (4.0,1.55);
  \draw[main!70!black, very thick] (t1) -- (t2);
  \node[font=\scriptsize, main!60!black] at (2.3,1.35) {S};
  \draw[third!55!black, dashed] (4.85,1.85) circle (0.85);
  \draw[main!70!black, very thick] (t2) arc (200:265:0.85);
  \coordinate (g) at (4.85,1.0);
  \draw[->, main!70!black, thick] (g) -- ++(0.58,-0.16);
  \fill[main!70!black] (g) circle (1.3pt) node[below right, font=\scriptsize]{$\mathbf{q}_f$};
  \node[font=\scriptsize, second!50!black] at (-0.85,0.85) {L};
  \node[font=\scriptsize, third!50!black] at (5.75,1.85) {L};
  \node[font=\scriptsize\bfseries, gray] at (2.4,-0.55) {Dubins LSL（只前进，无尖点）};
  % --- Reeds-Shepp with one cusp (bottom): forward arc then backward arc ---
  \begin{scope}[yshift=-2.7cm]
  \coordinate (rs) at (0,0);
  \draw[->, main!70!black, thick] (rs) -- ++(0.6,0);
  \fill[main!70!black] (rs) circle (1.3pt) node[below left, font=\scriptsize]{$\mathbf{q}_i$};
  % forward left arc
  \draw[main!70!black, very thick] (0,0) arc (-90:-10:0.95);
  \coordinate (cusp) at ($(0,0.95)+(-10:0.95)$);
  % cusp marker
  \fill[third!70!black] (cusp) circle (1.6pt);
  \node[font=\scriptsize, third!65!black, above right=0pt of cusp]{尖点 $\{\zeta\}^\flat$};
  % backward right arc (dashed-style = reverse), curving the other way down to goal
  \draw[main!55!black, very thick, densely dotted] (cusp) arc (110:30:0.95);
  \coordinate (rg) at ($(cusp)+(1.42,-0.62)$);
  \draw[<-, main!70!black, thick] (rg) -- ++(0.55,-0.2);
  \fill[main!70!black] (rg) circle (1.3pt) node[below right, font=\scriptsize]{$\mathbf{q}_f$};
  \node[font=\scriptsize, main!60!black] at (0.55,0.95) {$C^{+}$};
  \node[font=\scriptsize, main!50!black] at (2.55,0.7) {$C^{-}$（倒车）};
  \node[font=\scriptsize\bfseries, gray] at (2.4,-0.95) {Reeds--Shepp（前进弧 $\to$ 尖点换向 $\to$ 倒车弧）};
  \end{scope}
\end{tikzpicture}
\caption{曲率受限最短路两族。上：Dubins LSL，全程前进、由两段半径 $\rho$ 圆弧经公切线相连（CSC），无尖点。下：Reeds--Shepp 在中途出现\emph{尖点}（实心点，$v$ 过零换向），前一段前进弧 $C^+$、后一段倒车弧 $C^-$。尖点处朝向由 \cref{eq:nh-flat-recon} 的单侧极限连续补出（\cref{def:nh-cusp}）。}
\label{fig:nh-cusp}
\end{figure}

\subsection{与微分平坦、弧长参数化的关系}
\label{sec:nh-optpath-flat}

这两族曲线与本章前半（链式型 + 平坦）的接口，是它们能当采样规划 \texttt{Steer} 的关键。

\begin{insight}[弧长参数化下，曲率即唯一输入]\label{ins:nh-arclength}
取\emph{弧长} $s$ 为参数（$\mathrm{d}s=|v|\,\mathrm{d}t$），独轮车的笛卡尔路径满足 $x'^2+y'^2=1$（撇为 $\mathrm{d}/\mathrm{d}s$，自动归一），而 \cref{eq:nh-flat-recon} 退化为
\begin{equation}
  \theta=\operatorname{Atan2}(y',x'),\qquad \kappa(s)=\theta'(s)=x'y''-y'x'' ,
  \label{eq:nh-arclength-curv}
\end{equation}
即朝向是切向、\emph{曲率} $\kappa$ 是唯一的“几何输入”（时间律 $v(t)$ 可事后任配）。于是“前向速率恒定、转弯半径有下界 $\rho$”这条非完整 + 执行器约束，在弧长域里就是一句\emph{曲率上界} $|\kappa(s)|\le1/\rho$。Dubins/Reeds--Shepp 最短路据此可读成：在 $|\kappa|\le1/\rho$ 约束下、连接两端\emph{位置 + 切向}的\emph{最短弧长}曲线——其最优解必是 $\kappa\in\{-1/\rho,0,+1/\rho\}$ 的\emph{分段常曲率}拼接（bang--bang 在曲率上的体现，\cref{ins:nh-dubins-why}），正对应“最大曲率弧 + 直线段”。
\end{insight}

\begin{note}[平坦给可行性、最优给最短：两类 steering 的合流]\label{note:nh-steer-synthesis}
\cref{sec:nh-steer} 的多项式/平坦转向在平坦输出空间\emph{随意插值}，保证\emph{可行}但不保证\emph{最短}（且端点曲率不一定满足 $|\kappa|\le1/\rho$）；本节的 Dubins/Reeds--Shepp 给\emph{曲率受限下的最短}，但只覆盖“恒速、单一最小半径”的特定车模型。二者互补：采样规划器常以 Reeds--Shepp 解析路径作两点 \texttt{Steer} 与距离，再用平坦多项式或 \cref{sec:nh-track} 的跟踪控制器把它\emph{平滑化、时间参数化}并闭环执行。一句话——\emph{平坦管“连得上”，最优路径管“连得短”，跟踪管“跟得住”}。注意尖点处 $v=0$，恰落在 \cref{sec:nh-track} 跟踪控制器“失去持续激励”的奇异工况（\cref{tab:nh-track}），故含尖点路径须分段（以尖点断开）跟踪，或在尖点邻域切到 \cref{sec:nh-reg} 的镇定律。
\end{note}

% ════════════════════════════════════════════════════════════════
\section{可控性与可积性速查}
\label{sec:nh-quickguide}

本章与 \cref{ch:mm-wheeled} 反复用到“判约束可积否、判系统可达否”两套李括号判据。这里不重证（理论见 \cref{thm:cd-frobenius,sec:ctrl-controllability,ins:nh-twoface}），只把它们压成两段\emph{可照着算的流程}与一张\emph{四底盘速查表}，供动手时直接调用。两套判据用的是\emph{同一个}李括号，只是问向相反：Frobenius 问“括号\emph{落不落进}速度分布”（落进 $\Rightarrow$ 可积、完整），Chow 问“括号\emph{补不补满}全空间”（补满 $\Rightarrow$ 可达、可控）——这正是 \cref{ins:nh-twoface} 的“一体两面”。

\begin{algo}[Frobenius 可积判定：Pfaffian 约束是完整还是非完整]\label{alg:nh-frobenius}
判定一组 Pfaffian 速度约束\emph{能否积}成位形等式：算速度分布基两两的李括号，看有没有一个\emph{逸出}该分布——逸出即不可积（非完整）。
\end{algo}
\begin{codebox}[text]{Frobenius 可积判定（输入 A(q)q̇=0，k 行秩 k，位形维 n）}
Frobenius-Integrable(A, n)
 1  D <- ker A(q);  取基 {g_1, ..., g_m},  m = n - k     // 约束分布及其基
 2  for 每个基对 (i, j):
 3      B <- [g_i, g_j] = (dg_j/dq) g_i - (dg_i/dq) g_j  // 李括号
 4      if rank[ g_1 ... g_m  B ] > m:                    // B 逸出 span{g_*}
 5          return NONHOLONOMIC      // D 非对合 => 不可积（Thm Frobenius）
 6  return HOLONOMIC                 // D 对合 => 可积为 h(q)=0，运动困在子流形
\end{codebox}

\begin{algo}[Chow--Rashevsky / LARC 可达判定：无漂移系统能否到处去]\label{alg:nh-chow}
从输入向量场出发，\emph{反复添加}它们与已得分布的李括号，直到分布不再增长；若最终张满全空间则系统可达（STLC），所用的最高括号阶给出非完整度 $\kappa$。
\end{algo}
\begin{codebox}[text]{Chow / LARC 可达判定（无漂移系统，位形维 n）}
Chow-Reachable({g_1,...,g_m}, n)
 1  Δ_1 <- span{g_1, ..., g_m};   l <- 1
 2  while dim Δ_l < n  and  Δ_l != Δ_(l-1):              // 分布仍在长
 3      Δ_(l+1) <- Δ_l + span{ [g, h] : g  in  Δ_1, h  in  Δ_l }   // 加一阶李括号
 4      l <- l + 1
 5  if dim Δ_l = n （处处成立）:
 6      kappa <- l                       // 张满全空间所需的层数 = 非完整度
 7      return STLC-CONTROLLABLE, kappa  // LARC 成立（见 ins:nh-twoface）
 8  else:
 9      return NOT-CONTROLLABLE      // 可达集为 dim Δ_l 维子流形
\end{codebox}

\begin{insight}[读这两段流程的三个要点]\label{ins:nh-quickguide-read}
其一，\textbf{同源对照}：\cref{alg:nh-frobenius} 的“括号是否逸出 $\mathcal{D}$”与 \cref{alg:nh-chow} 的“括号是否补满 $\mathbb{R}^n$”用同一组括号，前者要它\emph{进}、后者要它\emph{出}。其二，\textbf{非完整度 $\kappa$ 的意义}：\cref{alg:nh-chow} 里滤链 $\Delta_1\subset\Delta_2\subset\cdots$ 张满全空间所需的层数即 $\kappa$，它量化“要凑出某些方向得做几阶机动”——$\kappa$ 越大规划越难（Ball--Box 度量代价 $\sim\varepsilon^{1/\kappa}$，\cref{ch:mm-wheeled}），与链式型 \cref{prop:nh-chained-ctrl} 的 $\kappa=n-1$ 一致。其三，\textbf{与线性判据的接续}：线性系统 $\operatorname{ad}_{\mathbf{f}}\mathbf{g}=-\mathbf{A}\mathbf{g}$，\cref{alg:nh-chow} 的滤链退化为可控性矩阵 $[\mathbf{B},\mathbf{A}\mathbf{B},\dots]$ 的列空间（\cref{sec:ctrl-controllability}），故本流程是线性可控性判据的非线性几何推广。
\end{insight}

\begin{table}[H]\centering\small
\caption{四类底盘的可控性 / 可积性速查（无漂移运动学模型）}
\label{tab:nh-chassis-rank}
\begin{tabular}{lcccccl}
\toprule
底盘 & 位形维 $n$ & 输入维 $m$ & 约束数 $k$ & 最高括号阶 & 非完整度 $\kappa$ & 可达性 \\
\midrule
独轮车 / 差速 & $3$ & $2$ & $1$ & $1$ & $2$ & STLC（$[\mathbf{g}_1,\mathbf{g}_2]$ 补横向） \\
自行车 / Ackermann & $4$ & $2$ & $2$ & $2$ & $3$ & STLC（化 $(2,4)$ 链式型） \\
单拖挂（$1$ 节挂车） & $4$ & $2$ & $2$ & $2$ & $3$ & STLC（链再延一级） \\
$N$ 节拖挂链 & $3{+}N$ & $2$ & $1{+}N$ & $1{+}N$ & $2{+}N$ & STLC（$\kappa$ 随挂车线性增） \\
\bottomrule
\end{tabular}
\end{table}

\begin{insight}[速查表怎么读：链越长越“硬”]\label{ins:nh-chassis-read}
四类底盘都\emph{可控}（STLC 列全是“是”），印证 \cref{ins:nh-essence}——非完整系统瞬时受限、整体可达。但“可控”不等于“好规控”：在\emph{正则位形}（无相邻两挂车互相垂直）下，非完整度 $\kappa$ 随拖挂节数\emph{线性增长}（独轮 $\kappa=2$、自行车/单挂 $\kappa=3$、独轮挂 $N$ 节车 $\kappa=2+N$，即 Jean 的“挂车数 $+2$”）。$\kappa$ 越大，凑出最后那个方向要做的李括号阶数越高、对应机动越“费事”，Ball--Box 代价越大、平坦输出的导数阶 $r$ 越高、链式型链越长（\cref{prop:nh-chained-ctrl}），规划与镇定都越难。（注：在\emph{奇异位形}——相邻挂车两两垂直——处 $\kappa$ 退化为 Fibonacci 数 $F_{N+3}$、随 $N$ 指数增，故“线性增”仅就一般位形而言；细节见 \cref{ch:mm-wheeled}。）这就是“多挂一节车，倒车入库就难一截”的数学解释。注意：所有这些底盘都\emph{非完整}（\cref{alg:nh-frobenius} 对每个都返回“非完整”，因横向约束的李括号必逸出速度分布），且都因 Brockett 障碍\emph{无法光滑镇定到一点}（\cref{thm:nh-brockett}）——可控性、可积性、可镇定性三者环环相扣，构成本章的逻辑闭环。
\end{insight}

\begin{pitfall}[速查表只对“纯运动学、无漂移”模型成立]\label{pit:nh-quickguide-scope}
\cref{tab:nh-chassis-rank} 与两段流程处理的是\emph{一阶、无漂移}模型 $\dot{\mathbf{q}}=\sum u_i\mathbf{g}_i$（输入归零即静止，\cref{eq:nh-affine}）。一旦把惯性、打滑、执行器动力学算进来，系统出现\emph{漂移项} $\mathbf{f}_0(\mathbf{q})$，可控性判据要换成更细的 STLC 充分条件（Sussmann 一般定理、好/坏括号配平），\cref{alg:nh-chow} 的“纯李括号补秩”不再充分。那属\cref{ch:underactuated} 与本部动力学诸章的二阶非完整 / 欠驱动范畴；本速查严格限定在轮式底盘的纯运动学层。
\end{pitfall}

% ════════════════════════════════════════════════════════════════
\section{可控性李代数的结构：自由李代数、增长向量与 Philip Hall 基}
\label{sec:nh-lie-structure}

\cref{alg:nh-chow} 的可达判定靠“反复添加李括号、看分布何时长满”，\cref{prop:nh-chained-ctrl} 也用重复括号 $\operatorname{ad}_{\boldsymbol{\gamma}_1}^k\boldsymbol{\gamma}_2$ 算非完整度。但两处都\emph{回避}了一个本该追问的问题：那些越来越长的迭代括号，到底\emph{有多少个互相独立}？若像 \cref{alg:nh-chow} 那样把所有 $[\mathbf{g}_i,[\mathbf{g}_j,\mathbf{g}_k]]$ 一股脑列出来，会被反对称 $[\mathbf{g},\mathbf{g}]=\mathbf 0$、$[\mathbf{g}_i,\mathbf{g}_j]=-[\mathbf{g}_j,\mathbf{g}_i]$ 和 Jacobi 恒等式\emph{大量重复计数}。本节补上这层代数骨架：先把可控性李代数及其\emph{滤过}讲清、给\emph{增长向量}一个干净定义（消除 \cref{thm:nh-sin-steer} 与 \cref{ins:nh-quickguide-read} 中对“增长向量”的名义提及），再用 Witt 公式数清自由李代数每层的真实维数，最后用 \textbf{Philip Hall 基}给出一组\emph{无冗余}的迭代括号——这正是把 \cref{alg:nh-chow} 从“会算一个例子”变成“能系统地、不重不漏地构造”的关键\cite{murray1994mathematical}。

\subsection{可控性李代数与它的滤过}
\label{sec:nh-ctrl-lie}

\begin{definition}[李代数与可控性李代数]\label{def:nh-ctrl-lie}
向量空间 $L$（实数域上）配一个双线性运算 $[\,\cdot\,,\cdot\,]:L\times L\to L$，若满足\emph{反对称} $[X,Y]=-[Y,X]$ 与 \emph{Jacobi 恒等式} $[X,[Y,Z]]+[Y,[Z,X]]+[Z,[X,Y]]=\mathbf 0$，则称 $L$ 为\emph{李代数}。$\mathbb{R}^n$ 上全体光滑向量场配李括号即一个（无穷维）李代数。给定无漂移系统 \cref{eq:nh-affine} 的输入向量场 $\mathbf{g}_1,\dots,\mathbf{g}_m$，包含它们的\emph{最小}李代数——即对李括号封闭的最小向量场集——记作
\begin{equation}
  \overline{\Delta}=\mathcal{L}(\mathbf{g}_1,\dots,\mathbf{g}_m),
  \label{eq:nh-ctrl-lie}
\end{equation}
称\emph{可控性李代数}（亦即分布 $\Delta=\operatorname{span}\{\mathbf{g}_i\}$ 的\emph{对合闭包}）。其元素由 $\mathbf{g}_i$ 的一切线性组合、它们的李括号、括号的括号……取遍得到。$\mathcal{L}(\mathbf{g}_1,\dots,\mathbf{g}_m)$ 在点 $\mathbf{q}$ 的\emph{秩}定义为 $\overline{\Delta}_{\mathbf{q}}$ 作为子空间的维数。
\end{definition}

要量化“凑出某个方向需要几阶括号”，把 $\overline{\Delta}$ 按括号长度分层：

\begin{definition}[滤过（filtration）]\label{def:nh-filtration}
令 $\Delta_1=\Delta=\operatorname{span}\{\mathbf{g}_1,\dots,\mathbf{g}_m\}$，递推定义
\begin{equation}
  \Delta_i=\Delta_{i-1}+[\Delta_1,\Delta_{i-1}],\qquad
  [\Delta_1,\Delta_{i-1}]\doteq\operatorname{span}\{[\mathbf{g},\mathbf{h}]:\mathbf{g}\in\Delta_1,\ \mathbf{h}\in\Delta_{i-1}\}.
  \label{eq:nh-filtration}
\end{equation}
则 $\Delta_1\subseteq\Delta_2\subseteq\cdots$，称与 $\Delta$ 相伴的\emph{滤过}。$\Delta_i$ 由输入向量场加上\emph{至多 $i-1$ 阶}李括号张成。若 $\operatorname{rank}\Delta_i(\mathbf{q})$ 在 $\mathbf{q}_0$ 邻域内不随 $\mathbf{q}$ 变，称该滤过在此\emph{正则}。
\end{definition}

\begin{proposition}[滤过的分级相容性：{$[\Delta_i,\Delta_j]\subseteq\Delta_{i+j}$}]\label{prop:nh-graded}
\cref{eq:nh-filtration} 的滤过满足 $[\Delta_i,\Delta_j]\subseteq\Delta_{i+j}$（约定 $\Delta_0=\{\mathbf 0\}$）。
\end{proposition}
\begin{proof}
对 $i$ 归纳。$i=1$ 即定义 \cref{eq:nh-filtration}：$[\Delta_1,\Delta_j]\subseteq\Delta_{j+1}$。设 $[\Delta_{i-1},\Delta_j]\subseteq\Delta_{i-1+j}$ 对一切 $j$ 成立。$\Delta_i=\Delta_{i-1}+[\Delta_1,\Delta_{i-1}]$，故只需证 $[[\Delta_1,\Delta_{i-1}],\Delta_j]\subseteq\Delta_{i+j}$。取 $X\in\Delta_1,Y\in\Delta_{i-1},Z\in\Delta_j$，由 Jacobi 恒等式
\[
  [[X,Y],Z]=[X,[Y,Z]]-[Y,[X,Z]].
\]
右边第一项：$[Y,Z]\in[\Delta_{i-1},\Delta_j]\subseteq\Delta_{i-1+j}$（归纳假设），再括 $X\in\Delta_1$ 得 $[X,[Y,Z]]\in[\Delta_1,\Delta_{i-1+j}]\subseteq\Delta_{i+j}$。第二项：$[X,Z]\in[\Delta_1,\Delta_j]\subseteq\Delta_{j+1}$，再与 $Y\in\Delta_{i-1}$ 括，$[Y,[X,Z]]\in[\Delta_{i-1},\Delta_{j+1}]\subseteq\Delta_{i+j}$（归纳假设）。两项皆在 $\Delta_{i+j}$，得证。
\end{proof}

\begin{definition}[非完整度、增长向量、相对增长向量]\label{def:nh-growth}
设滤过 $\{\Delta_i\}$ 正则。使 $\operatorname{rank}\Delta_{\kappa+1}=\operatorname{rank}\Delta_{\kappa}$ 的\emph{最小}整数 $\kappa$ 称分布的\emph{非完整度}（degree of nonholonomy）——此后滤过停止增长，$\Delta_{\kappa}=\overline{\Delta}$。定义\emph{增长向量} $\mathbf r=(r_1,\dots,r_{\kappa})$、\emph{相对增长向量} $\boldsymbol\sigma=(\sigma_1,\dots,\sigma_{\kappa})$：
\begin{equation}
  r_i=\operatorname{rank}\Delta_i,\qquad \sigma_i=r_i-r_{i-1}\ (r_0\doteq0).
  \label{eq:nh-growth}
\end{equation}
$r_i$ 记“用到 $i$ 阶括号时能张多少维”，$\sigma_i$ 记“第 $i$ 层\emph{新增}多少独立方向”。
\end{definition}

\noindent 至此 \cref{ins:nh-twoface,ins:nh-quickguide-read} 的可控性条件可一句话收紧：\textbf{LARC（Chow）成立 $\iff r_{\kappa}=\operatorname{rank}\overline{\Delta}=n$}。\cref{prop:nh-chained-ctrl} 的链式型恰是 $\sigma=(2,1,1,\dots,1)$、$\kappa=n-1$ 的“最瘦”滤过——每层只新增一个方向，故非完整度顶到 $n-1$。

\subsection{冗余从何来：Witt 公式数清每层的独立括号}
\label{sec:nh-witt}

\cref{eq:nh-filtration} 每升一层要把 $\Delta_1$ 与 $\Delta_{i-1}$ 的元素两两作括号，候选括号数随阶数爆炸；但反对称与 Jacobi 让其中大多数\emph{线性相关}。要数清第 $i$ 层最多能有几个独立括号，须先回答纯代数问题：\emph{$m$ 个生成元的自由李代数，长度恰为 $d$ 的齐次分量维数是多少？}

\begin{theorem}[Witt 公式：自由李代数齐次分量的维数\cite{murray1994mathematical}]\label{thm:nh-witt}
$m$ 个生成元上的自由李代数，长度恰为 $d$ 的齐次子空间维数为
\begin{equation}
  N(m,d)=\frac{1}{d}\sum_{e\mid d}\mu(e)\,m^{\,d/e},
  \label{eq:nh-witt}
\end{equation}
其中 $\sum_{e\mid d}$ 遍历 $d$ 的正因子，$\mu$ 为 Möbius 函数（$\mu(1)=1$；$\mu$ 在含 $r$ 个不同素因子的无平方数取 $(-1)^r$；非无平方数取 $0$）。该数也是“$m$ 色长度 $d$ 项链（Lyndon 词）”的个数。
\end{theorem}

\noindent 对两输入系统（$m=2$，独轮车、自行车、拖挂链全在此列），\cref{eq:nh-witt} 给 \cref{tab:nh-witt}：

\begin{table}[H]\centering\small
\caption{$m$ 生成元自由李代数各长度的独立括号数 $N(m,d)$（Witt 公式 \cref{eq:nh-witt}）}
\label{tab:nh-witt}
\begin{tabular}{lcccccc}
\toprule
长度 $d$ & $1$ & $2$ & $3$ & $4$ & $5$ & $6$ \\
\midrule
$N(2,d)$（两输入） & $2$ & $1$ & $2$ & $3$ & $6$ & $9$ \\
$N(3,d)$（三输入） & $3$ & $3$ & $8$ & $18$ & $48$ & $116$ \\
\bottomrule
\end{tabular}
\end{table}

\begin{insight}[“增长向量 $(m,m(m{+}1)/2)$”从何而来]\label{ins:nh-growth-firstorder}
\cref{thm:nh-sin-steer} 末尾对一般 $m$ 输入一阶系统写的增长向量 $(m,\,m(m{+}1)/2)$，现在有了精确来由：第一层 $r_1=\dim\Delta_1=m$；第二层新增的独立括号正是 $[\mathbf{g}_i,\mathbf{g}_j]\ (i<j)$，反对称把 $i=j$ 与 $i>j$ 全消去，剩 $\binom{m}{2}=m(m-1)/2$ 个——恰是 $N(m,2)=\tfrac1{2}(m^2-m)$（\cref{tab:nh-witt} 第二列）。于是若这些括号在该点都独立，$r_2=r_1+N(m,2)=m+\tfrac{m(m-1)}2=\tfrac{m(m+1)}2$。一般地 $\sigma_i\le N(m,i)$：相对增长被自由李代数的“代数容量”\emph{从上}卡死，等号当且仅当该层括号在该点处处独立（“自由”/无额外几何退化）。这把“规划难度随阶数涨”量化到了可数的代数极限上。
\end{insight}

\subsection{Philip Hall 基：一组无冗余的迭代括号}
\label{sec:nh-philiphall}

知道了每层\emph{有多少}独立括号，还需一套规则\emph{挑出}它们、不靠碰运气。Philip Hall 基就是这样一组规则：它把反对称与 Jacobi 的约束“编进”一个全序里，自动避开重复。

\begin{definition}[Lie 积的长度与幂零阶]\label{def:nh-bracket-length}
对生成元集 $\{\mathbf{g}_1,\dots,\mathbf{g}_m\}$，递归定义 \emph{Lie 积}（嵌套括号）的\emph{长度}
\begin{equation}
  l(\mathbf{g}_i)=1\ (i=1,\dots,m),\qquad l([A,B])=l(A)+l(B),
  \label{eq:nh-bracket-length}
\end{equation}
即 $l(A)$ 等于 $A$ 展开式里出现的生成元个数。若存在整数 $k$ 使一切长度 $>k$ 的 Lie 积皆为零，称该李代数\emph{幂零}（nilpotent），最小的 $k$ 为\emph{幂零阶}。
\end{definition}

\begin{definition}[Philip Hall 基\cite{murray1994mathematical}]\label{def:nh-philiphall}
$\mathcal{L}(\mathbf{g}_1,\dots,\mathbf{g}_m)$ 的一个 \emph{Philip Hall 基}是一个\emph{全序} Lie 积集 $H=\{B_i\}$，满足：
\begin{enumerate}
  \item $\mathbf{g}_i\in H$，$i=1,\dots,m$；
  \item 若 $l(B_i)<l(B_j)$ 则 $B_i<B_j$（按长度排序，等长内部任意定序）；
  \item $[B_i,B_j]\in H$ \emph{当且仅当}
  \begin{enumerate}
    \item $B_i,B_j\in H$ 且 $B_i<B_j$，且
    \item 或者 $B_j=\mathbf{g}_k$ 为某生成元，或者 $B_j=[B_l,B_r]$ 满足 $B_l,B_r\in H$ 且 $B_l\le B_i$。
  \end{enumerate}
\end{enumerate}
该集是 $\mathcal{L}(\mathbf{g}_1,\dots,\mathbf{g}_m)$ 的一组基（证明超本书范围，见\cite{murray1994mathematical} 及其所引 Bourbaki / Serre）。幂零阶 $k$ 的 Philip Hall 基可这样构造：列出长度 $\le k$ 的全部 Lie 积，逐一用上述条件筛除不合格者；实践中可只验条件 3。
\end{definition}

\begin{insight}[三条规则各挡掉哪种冗余]\label{ins:nh-phall-why}
条件 1--2 把生成元放在最前、按长度递增排，保证“先短后长”。条件 3 是去冗余的核心：要求 $[B_i,B_j]$ 只在 $B_i<B_j$ 时入选（用全序\emph{钉死}反对称——$[B_j,B_i]$ 与 $[B_i,B_i]$ 自动出局）；而 3(b) 的 $B_l\le B_i$（当 $B_j=[B_l,B_r]$）则\emph{钉死 Jacobi}——它规定嵌套括号必须“左因子单调不减”，从而在 Jacobi 恒等式联系的三个括号里只保留一个标准代表。两条约束合起来，使保留下来的括号个数恰好降到 \cref{thm:nh-witt} 的 $N(m,d)$，不多不少。
\end{insight}

\begin{example}[三生成元、幂零阶 $3$ 的 Philip Hall 基\cite{murray1994mathematical}]\label{ex:nh-phall3}
$\mathbf{g}_1,\mathbf{g}_2,\mathbf{g}_3$ 生成、幂零阶 $3$ 的李代数，其一组 Philip Hall 基为
\[
\begin{array}{llll}
\mathbf{g}_1 & \mathbf{g}_2 & \mathbf{g}_3 & \\[2pt]
[\mathbf{g}_1,\mathbf{g}_2] & [\mathbf{g}_1,\mathbf{g}_3] & [\mathbf{g}_2,\mathbf{g}_3] & \\[2pt]
[\mathbf{g}_1,[\mathbf{g}_1,\mathbf{g}_2]] & [\mathbf{g}_1,[\mathbf{g}_1,\mathbf{g}_3]] & [\mathbf{g}_2,[\mathbf{g}_1,\mathbf{g}_2]] & [\mathbf{g}_2,[\mathbf{g}_1,\mathbf{g}_3]] \\[2pt]
[\mathbf{g}_2,[\mathbf{g}_2,\mathbf{g}_3]] & [\mathbf{g}_3,[\mathbf{g}_1,\mathbf{g}_2]] & [\mathbf{g}_3,[\mathbf{g}_1,\mathbf{g}_3]] & [\mathbf{g}_3,[\mathbf{g}_2,\mathbf{g}_3]]
\end{array}
\]
共 $3+3+8=14$ 个，分别等于 $N(3,1)+N(3,2)+N(3,3)=3+3+8$（\cref{tab:nh-witt}）。\textbf{注意 $[\mathbf{g}_1,[\mathbf{g}_2,\mathbf{g}_3]]$ \emph{不}出现}：由 Jacobi
\[
  [\mathbf{g}_1,[\mathbf{g}_2,\mathbf{g}_3]]+[\mathbf{g}_2,[\mathbf{g}_3,\mathbf{g}_1]]+[\mathbf{g}_3,[\mathbf{g}_1,\mathbf{g}_2]]=\mathbf 0
  \ \Longrightarrow\
  [\mathbf{g}_1,[\mathbf{g}_2,\mathbf{g}_3]]=[\mathbf{g}_2,[\mathbf{g}_1,\mathbf{g}_3]]-[\mathbf{g}_3,[\mathbf{g}_1,\mathbf{g}_2]],
\]
右边两项已在基中，故它是冗余的——这正是 \cref{def:nh-philiphall} 条件 3(b)（$B_l\le B_i$）要挡掉的（$[\mathbf{g}_1,[\mathbf{g}_2,\mathbf{g}_3]]$ 的内括号左因子 $\mathbf{g}_2>\mathbf{g}_1=B_i$，违例）。
\end{example}

\begin{example}[两生成元、阶 $3$ 的 Philip Hall 基即链式型的括号链]\label{ex:nh-phall2}
$m=2$、幂零阶 $3$ 时，Philip Hall 基为
\[
  \underbrace{\mathbf{g}_1,\ \mathbf{g}_2}_{\text{长 }1},\quad
  \underbrace{[\mathbf{g}_1,\mathbf{g}_2]}_{\text{长 }2},\quad
  \underbrace{[\mathbf{g}_1,[\mathbf{g}_1,\mathbf{g}_2]],\ [\mathbf{g}_2,[\mathbf{g}_1,\mathbf{g}_2]]}_{\text{长 }3},
\]
共 $2+1+2=5=N(2,1)+N(2,2)+N(2,3)$。把它与 \cref{prop:nh-chained-ctrl} 的 $\operatorname{ad}_{\boldsymbol{\gamma}_1}^k\boldsymbol{\gamma}_2$ 链对照：链式型只动用了其中“最右嵌套”那一支 $\mathbf{g}_2,[\mathbf{g}_1,\mathbf{g}_2],[\mathbf{g}_1,[\mathbf{g}_1,\mathbf{g}_2]],\dots$（即 $\operatorname{ad}_{\mathbf{g}_1}^k\mathbf{g}_2$），每层只取一个——这正是链式型 $\sigma=(2,1,1,\dots)$ 的代数写照：它把自由李代数“压”到一条单链上，故 $\kappa=n-1$ 最大、规划最“硬”。
\end{example}

\subsection{增长向量算例与规划复杂度}
\label{sec:nh-growth-examples}

把滤过、增长向量与 Philip Hall 基用到本章的标准底盘上，得 \cref{tab:nh-growth}。它细化了 \cref{tab:nh-chassis-rank} 的“非完整度”一列：同样的 $n$，相对增长向量 $\boldsymbol\sigma$ 还能进一步区分系统“硬”在哪一层。

\begin{table}[H]\centering\small
\caption{标准非完整系统的增长向量（无漂移运动学模型，正则位形）}
\label{tab:nh-growth}
\begin{tabular}{lcccl}
\toprule
系统 & 位形维 $n$ & 增长向量 $\mathbf r$ & 相对增长 $\boldsymbol\sigma$ & 非完整度 $\kappa$ \\
\midrule
独轮车 / 差速 & $3$ & $(2,3)$ & $(2,1)$ & $2$ \\
跳跃机器人（飞行相） & $3$ & $(2,3)$ & $(2,1)$ & $2$ \\
滚动圆盘 & $4$ & $(2,3,4)$ & $(2,1,1)$ & $3$ \\
运动学汽车（自行车） & $4$ & $(2,3,4)$ & $(2,1,1)$ & $3$ \\
球面手指滚动 & $5$ & $(2,3,5)$ & $(2,1,2)$ & $3$ \\
$(2,n)$ 链式型 & $n$ & $(2,3,\dots,n)$ & $(2,1,\dots,1)$ & $n-1$ \\
\bottomrule
\end{tabular}
\end{table}

\begin{example}[运动学汽车的滤过与增长向量 $(2,3,4)$]\label{ex:nh-car-growth}
后轮驱动自行车 \cref{eq:nh-bicycle}（$\mathbf{q}=(x,y,\theta,\phi)$）的两输入向量场及逐层括号为
\[
  \mathbf{g}_1=\begin{bmatrix}\cos\theta\\\sin\theta\\\tfrac1\ell\tan\phi\\0\end{bmatrix},\
  \mathbf{g}_2=\begin{bmatrix}0\\0\\0\\1\end{bmatrix},\
  \mathbf{g}_3=[\mathbf{g}_1,\mathbf{g}_2]=\begin{bmatrix}0\\0\\-\tfrac1{\ell\cos^2\phi}\\0\end{bmatrix},\
  \mathbf{g}_4=[\mathbf{g}_1,\mathbf{g}_3]=\begin{bmatrix}-\tfrac{\sin\theta}{\ell\cos^2\phi}\\[2pt]\tfrac{\cos\theta}{\ell\cos^2\phi}\\0\\0\end{bmatrix}.
\]
$\{\mathbf{g}_1,\mathbf{g}_2,\mathbf{g}_3,\mathbf{g}_4\}$ 在 $\phi\neq\pm\pi/2$ 处线性无关，故 $\Delta_1=\operatorname{span}\{\mathbf{g}_1,\mathbf{g}_2\}$（$r_1=2$）、$\Delta_2$ 添 $\mathbf{g}_3$（$r_2=3$）、$\Delta_3$ 添 $\mathbf{g}_4$（$r_3=4=n$）：增长向量 $(2,3,4)$、相对增长 $(2,1,1)$、$\kappa=3=n-1$。它与 $(2,4)$ 链式型同型（\cref{eq:nh-chained}、\cref{tab:nh-growth} 末行 $n=4$），故 \cref{note:nh-chainable} 说“$n\le4$ 两输入系统总能化链式型”在此兑现。系统在 $\phi=\pm\pi/2$（车头垂直、$\mathbf{g}_1$ 失定义）处滤过退化，须避开。
\end{example}

\begin{insight}[相对增长向量量化“硬度”：同 $n$ 未必同难]\label{ins:nh-growth-hardness}
\cref{tab:nh-growth} 里运动学汽车与球面手指都是 $\kappa=3$，但相对增长不同：汽车 $(2,1,1)$、球面手指 $(2,1,2)$。$\boldsymbol\sigma$ 越“靠后才补满”、最高阶括号占比越大，机动越费事——这与 \cref{ins:nh-chassis-read} 的 Ball--Box 代价 $\sim\varepsilon^{1/\kappa}$ 同源：度量复杂度由滤过\emph{最深}那层主导。链式型 $\sigma=(2,1,\dots,1)$ 是\emph{每层只补一个}的极端（$\kappa=n-1$ 最大），故在所有 $n$ 维两输入系统里规划最“硬”；而球面手指能在第三层一次补两维（$\sigma_3=2$），$\kappa$ 反比链式型小、相对好规划。一句话：\textbf{增长向量是比单一非完整度 $\kappa$ 更细的“可规控难度指纹”}，Philip Hall 基则是把这枚指纹一层层算出来的工具。
\end{insight}

\begin{practice}
\begin{enumerate}
  \item 用 Möbius 函数验证 \cref{tab:nh-witt} 的 $N(2,4)=3$ 与 $N(2,5)=6$、$N(3,3)=8$；并解释为何 $N(m,1)=m$ 恒成立。
  \item 对独轮车 \cref{eq:nh-unicycle} 直接算 $[\mathbf{g}_1,\mathbf{g}_2]$，验证滤过 $r_1=2,r_2=3$，写出增长向量与 $\kappa$，并指出它对应 \cref{ex:nh-phall2} Philip Hall 基的哪几个元素。
  \item 按 \cref{def:nh-philiphall} 的三条规则，从 $\{\mathbf{g}_1,\mathbf{g}_2\}$ 出发构造幂零阶 $4$ 的 Philip Hall 基，验证元素个数为 $N(2,1)+\cdots+N(2,4)=2+1+2+3=8$。
  \item （综合）解释 \cref{ex:nh-phall3} 中为何 $[\mathbf{g}_3,[\mathbf{g}_1,\mathbf{g}_2]]$ \emph{在}基中、而 $[\mathbf{g}_1,[\mathbf{g}_2,\mathbf{g}_3]]$ \emph{不在}——用条件 3(b) 的 $B_l\le B_i$ 逐一核对。
\end{enumerate}
\end{practice}
